
\documentclass[a4paper,8pt]{extarticle} % extarticle allows to use font size of 8pt.

\usepackage[a4paper, top=1.6cm, bottom=2cm, left=1.6cm, right=1.6cm]{geometry} % Marge reduction.

%% Language specific package
\usepackage[french]{babel}
\frenchbsetup{StandardLists=true} % Necessary to use enumitem with babel/french.

%% Font and typing packages
\usepackage{fontspec}
\setmainfont[
	Ligatures=TeX,
	ItalicFont={Dancing Script},
	BoldItalicFont={Dancing Script}
	]{PT Serif} % default is Latin Modern
\newfontfamily\antiquefont[Ligatures=TeX]{Caslon Antique} % fancy font
\usepackage{microtype}			% Greatly improves general appearance of the text.
\usepackage{SIunits}			% Unit appearance.
\usepackage{xspace}				% Define commands that appear not to eat spaces.
\usepackage{ulem}				% To cross words out. Use \sout{}.

%% Array utilities
\usepackage{array}				% Additionnal options for arrays.
\usepackage{colortbl}			% Additionnal options for coloring arrays.
\usepackage[table]{xcolor}		% Auto alternate grey-white rows.
\usepackage[export]{adjustbox}		% Centered pics in tables

%% List utilities
\usepackage[inline]{enumitem}   % Display inline lists.
\usepackage{etoolbox}           % General utility. Good for lists for instance.
\usepackage{xparse}             % List utilities.
\usepackage{datatool}	% Handling alphabetical order.

%% Frames
\usepackage{framed}				% Boxes.
\usepackage[framemethod=TikZ]{mdframed}% For fancy frames.
\usepackage{tikz}				% For fancy frames.
\usepackage{wrapfig}			% Fancy insertion of pics in text.

%% Page utilities
\usepackage{multicol}			% Allows to divide a part of the page in multiple columns.
	
%% Others
\usepackage{keyval}             % Used to create maps of commands/labels/objects.
	\makeatletter                  % Mandatory for the usage of keyval.
\usepackage{xstring}            % String parsing, cutting, etc.
\usepackage{hyperref} % Links in PDF.


%%% Update of the dotfill command to always get dots

\newcommand{\predotfill}{\penalty0\hbox{}\nobreak}%


%%% Command to avoid typing \xspace when creating a new name macro

\newcommand{\newnamemacro}[2]{\newcommand{#1}{#2}} % \xspace removed for compatibility with alphabetical ordering

%%% Language specific stuff


\newcommand{\translationteam}{\item \og AEnoriel \fg \item \og Anglachel \fg \item \og Astadriel \fg \item \og Batcat \fg \item \og Bigfish \fg \item \og Eru \fg  \item \og Gandarin \fg \item \og Groumbahk \fg \item \og Iluvatar \fg \item \og Mammstein \fg \item \og Shlagrabak \fg \item et beaucoup d'autres...}

\hypersetup{
	pdfauthor={Équipe de traduction française de T9A},
	pdfsubject={Règles pour le jeu Batailles Fantastiques : Le 9\ieme{} Âge},
}

%%% Commands %%%

\ifdef{\isitanAB}{

\newcommand{\addtosortedlist}[1]{%
	\protected@edef\textarg{#1}%
	\protected@edef\textwithoutspaces{\expandafter\removespaces\expandafter{\textarg}}%
	\substitute\textwithoutspaces{É}{E}% Most used special characters of the language, and equivalent for alphabetical ordering
	\substitute\textwithoutspaces{È}{E}%
	\substitute\textwithoutspaces{Ê}{E}%
	\substitute\textwithoutspaces{é}{e}%
	\substitute\textwithoutspaces{è}{e}%
	\substitute\textwithoutspaces{ê}{e}%
	\substitute\textwithoutspaces{À}{A}%
	\substitute\textwithoutspaces{à}{a}%
	\substitute\textwithoutspaces{ù}{u}%
	\expandafter\sortitem\expandafter[\textwithoutspaces]{#1}%
}%

\newcommand{\pts}[1]{% First step is to remove spaces if there are some
	\def\numberwithoutspaces{\expandafter\removespaces\expandafter{#1}}%
	% Next step is getting rid of formatting if there are any (bold, color, ...)
	\pdfstringdef\cleannumber{\numberwithoutspaces}%
	% Now we can try if it is 1 or not
	\expandafter\ifstrequal\expandafter{\cleannumber}{1}{#1~\labels@point}{%
	\expandafter\ifstrequal\expandafter{\cleannumber}{0.5}{0,5~\labels@point}{%
	\expandafter\ifstrequal\expandafter{\cleannumber}{1.5}{1,5~\labels@point}{%
	#1~\labels@points}}}%
}

}{}

% Dark gods
\newcommand{\dchange}{Changement}
\newcommand{\dlust}{Luxure}
\newcommand{\pestilence}{Pestilence}
\newcommand{\wrath}{Courroux}
\newcommand{\truechaos}{Chaos Primordial}


% Nothing to edit here

\ifdef{\isitanAB}{

\newcommand{\alliancepts}[1]{
\ifsubstring{#1}{\free}{%
		\free{}%
	}{%
	\ifsubstring{#1}{\permodel}{%
		\splitatinf{#1}\myoption\myvalue%
		\pts{\myvalue}\permodel{}%
	}{%
	\pts{#1}
	}}
}

% You might wanna change the order of the gods - advanced user
\newcommand{\allianceoptions}[1]{%
	\defallianceoptions{#1}%
	\unitentryformat{\labels@allianceoptions\spacebeforecolon{}:}\newline
	\expandafter\ifblank\expandafter{\allianceoptions@introsentence}{}{\noindent\allianceoptions@introsentence{}\spacebeforecolon{}:}
	
	\expandafter\ifblank\expandafter{\allianceoptions@wrath}{
		\setlength{\columnseprule}{0.5pt}
		\renewcommand{\columnseprulecolor}{\color{black!30}}
		\vspace*{-0.2cm}\begin{multicols}{3}\raggedcolumns
		
			\begin{center}
			\noindent\dchange{}
			
			\noindent\alliancepts{\allianceoptions@change}
			\vspace*{-0.3cm}
			\end{center}
		
		\columnbreak
		
			\begin{center}
			\noindent\dlust{}
			
			\noindent\alliancepts{\allianceoptions@lust}
			\vspace*{-0.3cm}
			\end{center}
		
		\columnbreak

			\begin{center}
			\noindent\pestilence{}
			
			\noindent\alliancepts{\allianceoptions@pestilence}
			\vspace*{-0.3cm}
			\end{center}
			
		\end{multicols}
		\setlength{\columnseprule}{0pt}	
	}{
		\setlength{\columnseprule}{0.5pt}
		\renewcommand{\columnseprulecolor}{\color{black!30}}
		\vspace*{-0.2cm}\begin{multicols}{4}\raggedcolumns
		
			\begin{center}
			\noindent\dchange{}
			
			\noindent\alliancepts{\allianceoptions@change}
			\vspace*{-0.3cm}
			\end{center}
		
		\columnbreak

			\begin{center}
			\noindent\wrath{}
			
			\noindent\alliancepts{\allianceoptions@wrath}
			\vspace*{-0.3cm}
			\end{center}
		
		\columnbreak
		
			\begin{center}
			\noindent\dlust{}
			
			\noindent\alliancepts{\allianceoptions@lust}
			\vspace*{-0.3cm}
			\end{center}
		
		\columnbreak

			\begin{center}
			\noindent\pestilence{}
			
			\noindent\alliancepts{\allianceoptions@pestilence}
			\vspace*{-0.3cm}
			\end{center}
			
		\end{multicols}
		\setlength{\columnseprule}{0pt}
	}
}

}{}


%%% Labels %%%

% Profile

\newcommand{\labels@M}{M}
\newcommand{\labels@WS}{CC}
\newcommand{\labels@BS}{CT}
\newcommand{\labels@S}{F}
\newcommand{\labels@T}{E}
\newcommand{\labels@W}{PV}
\newcommand{\labels@I}{I}
\newcommand{\labels@A}{A}
\newcommand{\labels@Ld}{Cd}
\newcommand{\labels@Invocation}{Invocation} % For Vampire Covenant profiles
\newcommand{\labels@roundbase}{rond} % printed after XX mm for round bases

\newcommand{\Strength}{Force}

% Technical

\newcommand{\labels@range}{Portée}
\newcommand{\labels@point}{pt}
\newcommand{\labels@points}{pts}
\newcommand{\labels@only}{uniquement}
\newcommand{\labels@magic}{Magie}
\newcommand{\labels@pathsused}{Génère ses sorts dans la Discipline}
\newcommand{\labels@model}{figurine}
\newcommand{\labels@models}{figurines}
\newcommand{\labels@Singlemodel}{Figurine \textbf{seule}}

% Unit entry labels

\newcommand{\labels@basesize}{Socle}
\newcommand{\labels@trooptype}{Type de troupe}
\newcommand{\labels@specialrules}{Règles Spéciales}
\newcommand{\labels@alignment}{Allégeance}
\newcommand{\labels@alliance}{Allégeance}
\newcommand{\labels@allianceoptions}{Options d'Allégeance}
\newcommand{\labels@greenhiderace}{Race de Peaux Vertes}
\newcommand{\labels@equipment}{Équipement}
\newcommand{\labels@weapons}{Armes}
\newcommand{\labels@armour}{Armure}
\newcommand{\labels@options}{Options}
\newcommand{\labels@commandgroup}{État-Major}
\newcommand{\labels@charactermounts}{Montures de Personnages}
\newcommand{\labels@mounts}{Montures}
\newcommand{\labels@mount}{Monture}
\newcommand{\labels@specialequipment}{Équipement Spécial}

% Command groups

\newcommand{\labels@champion}{Champion}
\newcommand{\labels@standardbearer}{Porte-Étendard}
\newcommand{\labels@musician}{Musicien}
\newcommand{\labels@singlebannerallowance}{Une seule unité de ce type peut prendre une Bannière Magique}
\newcommand{\labels@condsinglebannerallowance}{Une seule unité de ce type peut prendre une Bannière Magique si}
\newcommand{\labels@bannerallowance}{Bannière Magique}
\newcommand{\labels@veteranstandardbearer}{Peut devenir Porte-Étendard Vétéran}
\newcommand{\labels@championallowance}{Arme Magique}

% Titles

\newcommand{\labels@armylist}{Liste des Troupes}
\newcommand{\labels@lords}{Seigneurs}
\newcommand{\labels@heroes}{Héros}
\newcommand{\labels@coreunits}{Unités de Base}
\newcommand{\labels@specialunits}{Unités Spéciales}
\newcommand{\labels@rareunits}{Unités Rares}
\newcommand{\labels@armywiderules}{Règles Communes de l'Armée}
\newcommand{\labels@armyspecialrules}{Règles Spéciales de l'Armée}
\newcommand{\labels@armoury}{Armurerie}
\newcommand{\labels@magicalitems}{Objets Magiques}
\newcommand{\labels@magicalweapons}{Armes Magiques}
\newcommand{\labels@magicalarmour}{Armures Magiques}
\newcommand{\labels@talismans}{Talismans}
\newcommand{\labels@enchanteditems}{Objets Enchantés}
\newcommand{\labels@arcaneitems}{Objets Cabalistiques}
\newcommand{\labels@magicalbanners}{Bannières Magiques}
\newcommand{\labels@quickrefsheet}{Fiche de Référence}
\newcommand{\labels@changelog}{Change Log}

\newcommand{\labels@lordsInitial}{S}
\newcommand{\labels@heroesInitial}{H}
\newcommand{\labels@coreunitsInitial}{B}
\newcommand{\labels@specialunitsInitial}{S}
\newcommand{\labels@rareunitsInitial}{R}
\newcommand{\labels@mountsInitial}{M}


% Titlepage

\newcommand{\labels@fantasybattles}{Batailles Fantastiques}
\newcommand{\labels@NinthAge}{Le 9\ieme{} Âge}
\newcommand{\labels@armyrules}{Règles de l'Armée}
\newcommand{\labels@frontpagecredits}{%
\labels@fantasybattles{} : \labels@NinthAge{} est un jeu créé et entretenu par la communauté qui met en scène des affrontements de figurines.\newline
Toutes les règles sont disponibles gratuitement sur le site suivant. Vos retours et suggestions sont les bienvenus : \url{http://www.the-ninth-age.com/}
}
\newcommand{\labels@license}{Copyright Creative Commons license : \url{the-ninth-age.com/license.html}}
\newcommand{\labels@tableofcontents}{Sommaire}
\newcommand{\labels@introduction}{%
\begin{center}\noindent{\Largerfontsize\textbf{Note des traducteurs}}\end{center}
\vspace{0.5cm}

Nous souhaitons remercier chaleureusement l'équipe à l'initiative du 9\ieme{} Âge pour leur motivation et leur travail continu pour faire vivre notre passion. Nous espérons que ce jeu saura développer les qualités pour plaire au plus grand nombre et réunir les joueurs, amateurs comme habitués des tournois, autour de règles amusantes et équilibrées, pour finalement s'imposer comme un standard du jeu de figurines. Une grande ambition qui ne pourra s'accomplir que \textbf{grâce à vous}, la communauté, via des retours constructifs, afin de modeler le jeu selon nos désirs. N'étant \textbf{en aucun cas à but lucratif}, le 9\ieme{} Âge part avec un avantage considérable. Les règles des éventuelles nouvelles sorties ne sont pas dictées par le besoin de vendre ces nouveautés. Vous pouvez choisir et acheter vos figurines où bon vous semble, il n'y a pas un unique revendeur toléré. Enfin, vous pouvez être assurés que tant que le 9\ieme{} Âge sera joué, vous disposerez d'un \textbf{support continu et régulier}, celui-ci étant offert par la communauté.

Concernant la traduction en elle-même, nous avons fait de notre mieux pour vous offrir une version de qualité, dont nous espérons qu'elle surpasse celle de la version originale ! Si vous constatez des coquilles, des erreurs, merci de nous les signaler en nous contactant sur le forum du 9\ieme{} Âge, dans le \textbf{sous-forum français} (\url{http://www.the-ninth-age.com/index.php?board/117-french/}). Vous y trouverez aussi les dernières mises à jour. \textbf{En cas de conflit d'interprétation avec la version originale, la version originale fait référence}.

\vspace{0.5cm}
Que ce jeu vous apporte d'innombrables heures de plaisir partagé !

\vspace{1cm}

\ifdef{\translationteam}{
	\begin{multicols}{3}
	\begin{itemize}
		\translationteam
	\end{itemize}
	\end{multicols}
}{}
}
\newcommand{\labels@rulechanges}{% blank ATM
}
\newcommand{\labels@latexcredit}{Document réalisé à l'aide de \LaTeX .}


%%% Technical commands

\newcommand{\only}[1]{(#1 uniquement)}
\newcommand{\free}{gratuit}
\newcommand{\upto}{jusqu'à}
\newcommand{\Upto}{Jusqu'à}
\newcommand{\unlimited}{pas de limite}
\newcommand{\permodel}{/fig.}
\newcommand{\listlastchoice}{ ou}
\newcommand{\notif}[1]{(pas #1)}
\newcommand{\wordand}{et}
\newcommand{\wordwith}{avec}
\newcommand{\ifNmodelsorless}[1]{(#1 figurines ou moins)}
\newcommand{\unitwith}{unité avec}
\newcommand{\From}{De} % From ... to ... models
\newcommand{\wordto}{à}
\newcommand{\wordAll}{Tous}
\newcommand{\spacebeforecolon}{ } % French put a space before colons
\newcommand{\minprice}{Coût min. :}
\newcommand{\mincostfor}{Coût min. pour}
\newcommand{\maxunitsize}{Taille max.}
\newcommand{\additionalfigscost}{Les figurines additionnelles coûtent}


%%% Special rules %%%

\newcommand{\ambush}{Embuscade}
\newcommand{\armourpiercing}[1]{Perforant\ifblank{#1}{}{ (#1)}}
\newcommand{\bodyguard}[1]{Garde du Corps\ifblank{#1}{}{ (#1)}}
\newcommand{\breathweapon}[1]{Attaque de Souffle\ifblank{#1}{}{ (#1)}}
\newcommand{\channel}{Canalisation}
\newcommand{\crushattack}{Attaque Écrasante}
\newcommand{\daemonicinstability}{Instabilité Démoniaque}
\newcommand{\devastatingcharge}{Charge Dévastatrice}
\newcommand{\distracting}{Distrayant}
\newcommand{\divineattacks}{Attaques Divines}
\newcommand{\engineer}{Ingénieur}
\newcommand{\ethereal}{Éthéré}
\newcommand{\fastcavalry}{Cavalerie Légère}
\newcommand{\fear}{Peur}
\newcommand{\fightinextrarank}{Combat avec un Rang Supplémentaire}
\newcommand{\fireborn}{Né du Feu}
\newcommand{\flamingattacks}{Attaques Enflammées}
\newcommand{\flammable}{Inflammable}
\newcommand{\frenzy}{Frénésie}
\newcommand{\fly}[1]{Vol\ifblank{#1}{}{ (#1)}}
\newcommand{\grindingattacks}[1]{Attaques de Broyage\ifblank{#1}{}{ (#1)}}
\newcommand{\hardtarget}{Camouflé}
\newcommand{\hatred}{Haine}
\newcommand{\hellfire}{Feu Démoniaque}
\newcommand{\hidden}{Caché}
\newcommand{\holyattacks}{Attaques Divines} % deprecated, still has to be filled. same as Divine Attacks.
\newcommand{\immunetopsychology}{Immunisé à la Psychologie}
\newcommand{\impacthits}[1]{Touches d'Impact\ifblank{#1}{}{ (#1)}}
\newcommand{\insignificant}{Insignifiant}
\newcommand{\largetarget}{Grande Cible}
\newcommand{\lethalstrike}{Coup Fatal}
\newcommand{\lightningattacks}{Attaques Foudroyantes}
\newcommand{\lightningreflexes}{Réflexes Foudroyants}
\newcommand{\lighttroops}{Troupe Légère}
\newcommand{\magicresistance}[1]{Résistance à la Magie\ifblank{#1}{}{ (#1)}}
\newcommand{\magicalattacks}{Attaques Magiques}
\newcommand{\metalshifting}{Fusion du Métal}
\newcommand{\moveorfire}{Mouvement ou Tir}
\newcommand{\multipleshots}[1]{Tirs Multiples\ifblank{#1}{}{ (#1)}}
\newcommand{\multiplewounds}[2]{Blessures Multiples\ifblank{#1}{}{ (#1\ifblank{#2}{)}{, #2)}}}
\newcommand{\notaleader}{Pas un Meneur}
\newcommand{\otherworldly}{D'Outre-Monde}
\newcommand{\pathmaster}[1]{Maître de la Voie\ifblank{#1}{}{ (#1)}}
\newcommand{\poisonedattacks}{Attaques Empoisonnées}
\newcommand{\quicktofire}{Tir Rapide}
\newcommand{\randommovement}[1]{Mouvement Aléatoire\ifblank{#1}{}{ (#1)}}
\newcommand{\randomattacks}[1]{Attaques Aléatoires\ifblank{#1}{}{ (#1)}}
\newcommand{\regeneration}[1]{Régénération\ifblank{#1}{}{ (#1+)}}
\newcommand{\reload}{Rechargez !}
\newcommand{\requirestwohands}{Arme à deux Mains}
\newcommand{\scythes}{Faux}
\newcommand{\scout}{Éclaireur}
\newcommand{\scouts}{Éclaireurs}
\newcommand{\stomp}[1]{Piétinement\ifblank{#1}{}{ (#1)}}
\newcommand{\strider}[1]{Guide\ifblank{#1}{}{ (#1)}}
\newcommand{\stubborn}{Tenace}
\newcommand{\stupidity}{Stupidité}
\newcommand{\skirmisher}{Tirailleur}
\newcommand{\skirmishers}{Tirailleurs}
\newcommand{\sweepingattack}{Attaque au Passage}
\newcommand{\swiftstride}{Course Rapide}
\newcommand{\thunderouscharge}{Charge Tonitruante}
\newcommand{\terror}{Terreur}
\newcommand{\toxicattacks}{Attaques Toxiques}
\newcommand{\unbreakable}{Indémoralisable}
\newcommand{\undead}{Mort-Vivant}
\newcommand{\unstable}{Instable}
\newcommand{\unwieldy}{Encombrant}
\newcommand{\vanguard}{Avant-Garde}
\newcommand{\volleyfire}{Tir de Volée}
\newcommand{\warplatform}{Plateforme de Guerre}
\newcommand{\wardsave}[1]{Sauvegarde Invulnérable\ifblank{#1}{}{ (#1+)}}
\newcommand{\weaponmaster}{Maître d'Ar\-mes}
\newcommand{\wizardconclave}[1]{Conclave de Sorciers\ifblank{#1}{}{ (#1)}}


%%% Magic %%%


% General

\newcommand{\Pathof}{Voie}

\newcommand{\battle}{Commune}

\newcommand{\anyofthebattlemagic}{dans n'importe laquelle des Voies Communes}
\newcommand{\ONLYanyofthebattlemagic}{Commune de votre choix}

\newcommand{\magiclevel}[1]{\ifnumcomp{#1}{<}{3}{Apprenti Magicien}{Maître Magicien} Niveau #1}
\newcommand{\Level}{Niveau}

\newcommand{\wizard}{Magicien}
\newcommand{\wizards}{Magiciens}

\newcommand{\learnedspell}{Sort Appris}
\newcommand{\learnedspells}{Sorts Appris}
\newcommand{\attributespell}{Attribut de la Voie}
\newcommand{\attributespells}{Attributs de la Voie}
\newcommand{\attributespellnumber}{A}
\newcommand{\traitspell}{Sort Caractéristique}
\newcommand{\traitspells}{Sorts Caractéristiques}
\newcommand{\traitspellnumber}{C}


\newcommand{\boundspell}[1]{Objet de Sort\ifblank{#1}{}{, Puissance #1}}
\newcommand{\boundspells}[1]{Objets de Sort\ifblank{#1}{}{, Puissance #1}}

% Casting Vocabulary

\newcommand{\lostfocus}{Perte de Concentration}
\newcommand{\miscast}{Fiasco}
\newcommand{\miscasts}{Fiascos}
\newcommand{\overwhelmingpower}{Pouvoir Irrésistible}

\newcommand{\breachintheveil}{Brèche dans le Voile}
\newcommand{\catastrophicdetonation}{Explosion Catastrophique}
\newcommand{\witchfire}{Feu de Sorcières}
\newcommand{\sorcerousbacklash}{Contrecoup Magique}
\newcommand{\amnesia}{Amnésie}

% Spell Types

\newcommand{\augment}{Amélioration}
\newcommand{\hex}{Malédiction}
\newcommand{\universal}{Universel}
\newcommand{\missile}{Projectile}
\newcommand{\damage}{Dégâts}
\newcommand{\direct}{Direct}
\newcommand{\focused}{Focalisé}
\newcommand{\vortex}{Vortex}
\newcommand{\ground}{Marqueur}
\newcommand{\linetemplate}{Gabarit de Ligne}
\newcommand{\specialTYPE}{Spécial}
\newcommand{\aura}{Aura}
\newcommand{\castersunit}{Unité du Lanceur}
\newcommand{\caster}{Lanceur}

\newcommand{\template}{Gabarit}

% Spell Durations

\newcommand{\lastsoneturn}{Dure un Tour}
\newcommand{\instant}{Immédiat}
\newcommand{\permanent}{Permanent}
\newcommand{\remainsinplay}{Reste en Jeu}


% Battle Magic

\newcommand{\alchemy}{de l'Alchimie}
\newcommand{\alchemyattribute}{Édit de Fer}
\newcommand{\alchemysignature}{Métal Fondu}
\newcommand{\alchemyspellone}{Lames Enchantées}
\newcommand{\alchemyspelltwo}{Corrosion Rampante}
\newcommand{\alchemyspellthree}{Manteau de Vif-Argent}
\newcommand{\alchemyspellfour}{Pieu d'Argent}
\newcommand{\alchemyspellfive}{Fléau de l'Acier}
\newcommand{\alchemyspellsix}{Transmutation en Or}

\newcommand{\death}{de la Mort}
\newcommand{\deathattribute}{Nuage de Désespoir}
\newcommand{\deathsignature}{Le Baiser de la Faucheuse}
\newcommand{\deathspellone}{Malédiction du Mortel}
\newcommand{\deathspelltwo}{Esprits Dévorants}
\newcommand{\deathspellthree}{Sangsue Psychique}
\newcommand{\deathspellfour}{Moisson d’Âmes}
\newcommand{\deathspellfive}{L’Abîme aussi te Regarde...}
\newcommand{\deathspellsix}{Maelström d’Âmes}

\newcommand{\fire}{du Feu}
\newcommand{\fireattribute}{Feu Déchaîné}
\newcommand{\firesignature}{Boule de Feu}
\newcommand{\firespellone}{Cascade Ardente}
\newcommand{\firespelltwo}{Épées Flamboyantes}
\newcommand{\firespellthree}{Jet de Flammes}
\newcommand{\firespellfour}{Traits Enflammés}
\newcommand{\firespellfive}{Remparts Incandescents}
\newcommand{\firespellsix}{Souffler sur les Braises}

\newcommand{\heavens}{des Cieux}
\newcommand{\heavensattribute}{Second Sceau}
\newcommand{\heavenssignature}{Aquilon}
\newcommand{\heavensspellone}{Bourrasque}
\newcommand{\heavensspelltwo}{Choc Foudroyant}
\newcommand{\heavensspellthree}{Conjonction Astrale}
\newcommand{\heavensspellfour}{Fléau du Ponant}
\newcommand{\heavensspellfive}{Déluge d'Éclairs}
\newcommand{\heavensspellsix}{Appel de la Comète}

\newcommand{\light}{de la Lumière}
\newcommand{\lightattribute}{Lumière Gardienne}
\newcommand{\lightsignature}{Éclat Brûlant}
\newcommand{\lightspellone}{Bouclier Protecteur}
\newcommand{\lightspelltwo}{Étincelle de Courage}
\newcommand{\lightspellthree}{Vitesse Fulgurante}
\newcommand{\lightspellfour}{Toile Scintillante}
\newcommand{\lightspellfive}{Distorsion Temporelle}
\newcommand{\lightspellsix}{Bannissement Divin}

\newcommand{\nature}{de la Nature}
\newcommand{\natureattribute}{Souffle de Vie}
\newcommand{\naturesignature}{Eaux Vivifiantes}
\newcommand{\naturespellone}{Maître de la Terre}
\newcommand{\naturespelltwo}{Le Trône de Chêne}
\newcommand{\naturespellthree}{Esprits des Bois}
\newcommand{\naturespellfour}{Croissance Estivale}
\newcommand{\naturespellfive}{Peau Rocailleuse}
\newcommand{\naturespellsix}{Créatures Souterraines}

\newcommand{\shadows}{des Ombres}
\newcommand{\shadowsattribute}{Course Parmi les Ombres}
\newcommand{\shadowssignature}{Miasmes Obscurs}
\newcommand{\shadowsspellone}{Orbe de Noirceur}
\newcommand{\shadowsspelltwo}{Partir en Fumée}
\newcommand{\shadowsspellthree}{Expérience de Mort Imminente}
\newcommand{\shadowsspellfour}{Char Vaporeux}
\newcommand{\shadowsspellfive}{Ombres Dévorantes}
\newcommand{\shadowsspellsix}{Scalpel Psychique}

\newcommand{\wilderness}{de la Sauvagerie}
\newcommand{\wildernessattribute}{La Chasse Sauvage}
\newcommand{\wildernesssignature}{La Bête qui Sommeille}
\newcommand{\wildernessspellone}{Essaim d’Insectes}
\newcommand{\wildernessspelltwo}{Rage Intérieure}
\newcommand{\wildernessspellthree}{Pieu de Rougebois}
\newcommand{\wildernessspellfour}{Calamité des Bois Sauvages}
\newcommand{\wildernessspellfive}{Tempête Furieuse}
\newcommand{\wildernessspellsix}{Métamorphose
Monstrueuse}

\newcommand{\eightpaths}{Octuple}



% Army Specific Magic

\newcommand{\butchery}{de la Boucherie}
\newcommand{\butcheryattribute}{Sang de Kholag}
\newcommand{\butcherysignature}{Briseur de Dents}
\newcommand{\butcheryspellone}{Buveur de Moelle}
\newcommand{\butcheryspelltwo}{Festin de Tripaille}
\newcommand{\butcheryspellthree}{Concasseur d’Os}
\newcommand{\butcheryspellfour}{Gobeur de Cervelle}
\newcommand{\butcheryspellfive}{Cœur de Troll}
\newcommand{\butcheryspellsix}{Gosier de Géant}

\newcommand{\change}{du Changement}
\newcommand{\changeattribute}{Vent du Changement}
\newcommand{\changesignature}{Feu Azur}
\newcommand{\changespellone}{Feu Rose}
\newcommand{\changespelltwo}{Vague du Changement}
\newcommand{\changespellthree}{Secrets Volés}
\newcommand{\changespellfour}{Règne de la Confusion}
\newcommand{\changespellfive}{Inéluctable Trahison}
\newcommand{\changespellsix}{Portail Éternel}

\newcommand{\thebiggreengods}{des Grands Dieux Verts}
\newcommand{\thebiggreengodsattribute}{Chopez-les !}
\newcommand{\thebiggreengodssignature}{L'Heure de la Raclée}
\newcommand{\thebiggreengodsspellone}{Coup de Boule}
\newcommand{\thebiggreengodsspelltwo}{Poings Bastonneurs}
\newcommand{\thebiggreengodsspellthree}{Même Pas Mal !}
\newcommand{\thebiggreengodsspellfour}{Grande Main Verte}
\newcommand{\thebiggreengodsspellfive}{Boum !}
\newcommand{\thebiggreengodsspellsix}{Le Gros Piétinement}

\newcommand{\thelittlegreengods}{des Petits Dieux Verts}
\newcommand{\thelittlegreengodsattribute}{Fourbe Larcin}
\newcommand{\thelittlegreengodssignature}{Œil Mauvais}
\newcommand{\thelittlegreengodsspellone}{Taillades Sournoises}
\newcommand{\thelittlegreengodsspelltwo}{Bénédiction de la Mère-Araignée}
\newcommand{\thelittlegreengodsspellthree}{Ça Démange ?}
\newcommand{\thelittlegreengodsspellfour}{Chut ! Pas un Bruit...}
\newcommand{\thelittlegreengodsspellfive}{J’vous Arrange Ça}
\newcommand{\thelittlegreengodsspellsix}{Malédiction de la Lune Verte}

\newcommand{\blackmagic}{de la Magie Noire}
\newcommand{\blackmagicattribute}{Soif d’Âmes}
\newcommand{\blackmagicsignature}{Furie de Moraec}
\newcommand{\blackmagicspellone}{Rafale Glaciale}
\newcommand{\blackmagicspelltwo}{Tourbillon de Lames}
\newcommand{\blackmagicspellthree}{Agonie Paralysante}
\newcommand{\blackmagicspellfour}{Marque de la Peur}
\newcommand{\blackmagicspellfive}{Trait d’Énergie Noire}
\newcommand{\blackmagicspellsix}{Terreur Noire}

\newcommand{\disease}{de la Maladie}
\newcommand{\diseaseattribute}{Bénédiction Nécrotique}
\newcommand{\diseasesignature}{Relents de Pestilence}
\newcommand{\diseasespellone}{Haleine Corruptrice}
\newcommand{\diseasespelltwo}{Toucher Putréfiant}
\newcommand{\diseasespellthree}{Excroissance Adipeuse}
\newcommand{\diseasespellfour}{Purge Parasitaire}
\newcommand{\diseasespellfive}{Malédiction du Lépreux}
\newcommand{\diseasespellsix}{Tourbillon Fétide}

\newcommand{\lust}{de la Luxure}
\newcommand{\lustattribute}{Masochisme}
\newcommand{\lustsignature}{Flagellation Démoniaque}
\newcommand{\lustspellone}{Grâce Hypnotique}
\newcommand{\lustspelltwo}{Valse Irrésistible}
\newcommand{\lustspellthree}{Hystérie}
\newcommand{\lustspellfour}{Fantasmagorie}
\newcommand{\lustspellfive}{Déchirement psychique}
\newcommand{\lustspellsix}{Chœur Dissonant}

\newcommand{\necromancy}{de la Nécromancie}
\newcommand{\necromancyattribute}{Tromper la Faucheuse}
\newcommand{\necromancysignature}{Adjuration des Morts}
\newcommand{\necromancyspellone}{Parodie de Vie}
\newcommand{\necromancyspelltwo}{Convocation Profanatoire}
\newcommand{\necromancyspellthree}{Sarabande Macabre}
\newcommand{\necromancyspellfour}{Regard de Setesh}
\newcommand{\necromancyspellfive}{Vol de Jeunesse}
\newcommand{\necromancyspellsix}{Malédiction des Morts}

\newcommand{\ruin}{de la Ruine}
\newcommand{\ruinattribute}{Hordes Sans Fin}
\newcommand{\ruinsignature}{Éclair Noir}
\newcommand{\ruinspellone}{Nourrissons-les...}
\newcommand{\ruinspelltwo}{Souiller le Sol}
\newcommand{\ruinspellthree}{La Faim}
\newcommand{\ruinspellfour}{Appel de la Tempête}
\newcommand{\ruinspellfive}{Rupture Sismique}
\newcommand{\ruinspellsix}{Pour Qui Sonne le Glas}

\newcommand{\forge}{de la Forge}
\newcommand{\forgeattribute}{Fournaise Haineuse}
\newcommand{\forgesignature}{Bouclier de Sombrefeu}
\newcommand{\forgespellone}{Rage Incendiaire}
\newcommand{\forgespelltwo}{Subjugation}
\newcommand{\forgespellthree}{Souffle de Haine}
\newcommand{\forgespellfour}{Anathème de Noirceur}
\newcommand{\forgespellfive}{Cendres Asphyxiantes}
\newcommand{\forgespellsix}{Flammes de la Forge}

\newcommand{\sands}{des Sables}
\newcommand{\sandsattribute}{Les Morts sans Repos}
\newcommand{\sandssignature}{Sirocco}
\newcommand{\sandsspellone}{Lames Maudites}
\newcommand{\sandsspelltwo}{Dessiccation Mortelle}
\newcommand{\sandsspellthree}{Frappes Vengeresses}
\newcommand{\sandsspellfour}{Jugement Divin}
\newcommand{\sandsspellfive}{Sables Mouvants}
\newcommand{\sandsspellsix}{Écho des Gloires
Passées}

\newcommand{\whitemagic}{de la Magie Blanche}
\newcommand{\whitemagicattribute}{Bouclier des Anciens}
\newcommand{\whitemagicsignature}{Traits de Lumière}
\newcommand{\whitemagicspellone}{Résurrection du Phénix}
\newcommand{\whitemagicspelltwo}{Volonté Inspirante}
\newcommand{\whitemagicspellthree}{Sentier Secret}
\newcommand{\whitemagicspellfour}{Bénédiction d’Amhar}
\newcommand{\whitemagicspellfive}{Fusion d’Artefact}
\newcommand{\whitemagicspellsix}{Cataclysme}

% Paths Initials

\newcommand{\alchemyInitials}{A}
\newcommand{\deathInitials}{M}
\newcommand{\fireInitials}{F}
\newcommand{\heavensInitials}{C}
\newcommand{\lightInitials}{L}
\newcommand{\natureInitials}{N}
\newcommand{\shadowsInitials}{O}
\newcommand{\wildernessInitials}{S}

\newcommand{\eightfoldInitials}{8}

\newcommand{\whitemagicInitials}{MB}
\newcommand{\blackmagicInitials}{MN}
\newcommand{\necromancyInitials}{N}
\newcommand{\sandsInitials}{S}
\newcommand{\forgeInitials}{F}
\newcommand{\biggreengodsInitials}{GDV}
\newcommand{\littlegreengodsInitials}{PDV}
\newcommand{\butcheryInitials}{B}
\newcommand{\ruinInitials}{R}
\newcommand{\diseaseInitials}{M}
\newcommand{\lustInitials}{L}
\newcommand{\changeInitials}{C}


%%% Other rules %%%

% Troop types rules

\newcommand{\combinedprofile}{Profil Combiné}
\newcommand{\cavalrysupport}{Soutien de Cavalerie}
\newcommand{\monstrousranks}{Rangs Monstrueux}
\newcommand{\monstroussupport}{Soutien Monstrueux}
\newcommand{\monsterranks}{Rang de Monstre}

\newcommand{\armoursave}{Sauvegarde d'Armure}
\newcommand{\frontrank}{Au Premier Rang}
\newcommand{\hardcover}{Couvert Lourd}
\newcommand{\holdyourground}{Tenez les Rangs}
\newcommand{\inspiringpresence}{Présence Charismatique}
\newcommand{\lightcover}{Couvert Léger}
\newcommand{\ordnance}{Artillerie}
\newcommand{\parry}{Parade}
\newcommand{\raisewounds}{Ressusciter des Figurines}
\newcommand{\recoverwounds}{Récupérer des PVs}
\newcommand{\rnf}{ordinaires}
\newcommand{\general}{Général}
\newcommand{\bsb}{Porteur de la Grande Bannière}
\newcommand{\cannotmarch}{Pas de Marche Forcée}
\newcommand{\veteranstandardbearer}{Porte-Étendard Vétéran}
\newcommand{\swirlingmelee}{Mêlée Tourbillonnante}
\newcommand{\scoringunit}{Unité de Capture}
\newcommand{\scoringunits}{Unités de Capture}


%%% Equipment %%%

\newcommand{\hw}{Arme de Base}
\newcommand{\pw}{Paire d'Armes}
\newcommand{\spear}{Lance}
\newcommand{\halberd}{Hallebarde}
\newcommand{\gw}{Arme Lourde}
\newcommand{\lance}{Lance de Cavalerie}
\newcommand{\lightlance}{Lance Légère}
\newcommand{\flail}{Fléau}

\newcommand{\throwingweapons}{Armes de Jet}
\newcommand{\shortbow}{Arc Court}
\newcommand{\bow}{Arc}
\newcommand{\longbow}{Arc Long}
\newcommand{\handgun}{Arquebuse}
\newcommand{\crossbow}{Arbalète}
\newcommand{\pistol}{Pistolet}
\newcommand{\braceofpistols}{Paire de Pistolets}	

\newcommand{\innatedefence}[1]{Protection Innée\ifblank{#1}{}{~(#1+)}}
\newcommand{\mountsprotection}[1]{Protection de Monture\ifblank{#1}{}{~(#1+)}}
\newcommand{\la}{Armure Légère}
\newcommand{\ha}{Armure Lourde}
\newcommand{\platearmour}{Armure de Plates}
\newcommand{\shield}{Bouclier}
\newcommand{\barding}{Caparaçon}

\newcommand{\cannon}{Canon}
\newcommand{\cannons}{Canons}
\newcommand{\catapult}{Catapulte}
\newcommand{\catapults}{Catapultes}
\newcommand{\volleygun}{Batterie de Tir}
\newcommand{\boltthrower}{Baliste}
\newcommand{\flamethrower}{Lance-Flammes}
\newcommand{\artilleryweapon}{Arme d'Artillerie}


%%% Troop types %%%

\newcommand{\characters}{Personnages}
\newcommand{\infantry}{Infanterie}
\newcommand{\monstrousinfantry}{Infanterie Monstrueuse}
\newcommand{\cavalry}{Cavalerie}
\newcommand{\monstrouscavalry}{Cavalerie Monstrueuse}
\newcommand{\swarm}{Nuée}
\newcommand{\swarms}{Nuées}
\newcommand{\warbeast}{Bête de Guerre}
\newcommand{\warbeasts}{Bêtes de Guerre}
\newcommand{\monster}{Monstre}
\newcommand{\monsters}{Monstres}
\newcommand{\monstrousbeast}{Bête Monstrueuse}
\newcommand{\monstrousbeasts}{Bêtes Monstrueuses}
\newcommand{\chariot}{Char}
\newcommand{\chariots}{Chars}
\newcommand{\riddenmonster}{Monstre Monté}
\newcommand{\riddenmonsters}{Monstres Montés}
\newcommand{\warmachine}{Machine de Guerre}
\newcommand{\warmachines}{Machines de Guerre}


%%% Terrain %%%

\newcommand{\water}{Eaux Peu Profondes}
\newcommand{\forest}{Forêt}
\newcommand{\impassableterrain}{Terrain Infranchissable}


%%% Profile wording

\newcommand{\oneperarmy}{Un par Armée}
\newcommand{\oneofakind}{Uni\-que}
\newcommand{\zerotoXchoice}[1]{0-#1 Choix}
\newcommand{\onechoiceonlyNOC}{(un seul choix)}
\newcommand{\onfootonly}{(à pied uniquement)}
\newcommand{\closecombatonly}{seulement au Corps à Corps}
\newcommand{\Xmodelsorless}[1]{(max. #1 figurines)}
\newcommand{\magicalitemsallowance}{Objets Magiques}
\newcommand{\magicalweaponallowance}{Arme Magique}
\newcommand{\notmagicalarmour}{(pas d'Armure Magique)}
\newcommand{\weapononechoice}{\optionschoice{Arme \onechoiceonlyNOC{} :}}
\newcommand{\weaponschoice}{\optionschoice{Armes :}}
\newcommand{\shootingweapononechoice}{\optionschoice{Arme de Tir \onechoiceonlyNOC{} :}}
\newcommand{\combatweapononechoice}{\optionschoice{Arme de Corps à Corps \onechoiceonlyNOC{} :}}
\newcommand{\combatweapononechoiceTWOCOL}{\optionschoiceTWOCOL{Arme de Corps à Corps \onechoiceonlyNOC{} :}}
\newcommand{\armouronechoice}{\optionschoice{Armure \onechoiceonlyNOC{} :}}
\newcommand{\magiclevelchoice}{\optionschoice{Magie \onechoiceonlyNOC{} :}}
\newcommand{\mustbecomeoneofthefollowing}{\optionschoice{\textbf{Doit} devenir au choix :}}
\newcommand{\mustbecomeoneofthefollowingNOC}{Doit devenir au choix :}
\newcommand{\musttakeoneormoreofthefollowing}{\optionschoice{\textbf{Doit} prendre au moins un choix :}}
\newcommand{\musttakeoneofthefollowing}{\optionschoice{\textbf{Doit} prendre un et un seul choix :}}
\newcommand{\musttakeoneofthefollowingNOC}{Doit choisir entre :}
\newcommand{\uptotwoofthefollowing}{\optionschoice{Jusqu'à deux choix :}}
\newcommand{\uptotwoofthefollowingTWOCOL}{\optionschoiceTWOCOL{Jusqu'à deux choix :}}

\newcommand{\onechoiceonly}{\optionschoice{Un seul choix :}}
\newcommand{\onechoiceonlyTWOCOL}{\optionschoiceTWOCOL{Un seul choix :}}

\newcommand{\maytake}{Peut prendre}




%%% Orcs N Goblins debug, let it as it is

\newcommand{\pershadygit}{debug}
\newcommand{\permadgit}{debug}

%%% Dwarven Holds debug, let it as it is

\newcommand{\perrune}{debug}


%%% Commands to handle strings, better than xstring to handle commands inside the strings %%%

\newcommand{\substitute}[3]{%
  \protected@edef\sub@temp{#1}%
  \saveexpandmode
  \expandarg\StrSubstitute{\sub@temp}{#2}{#3}[#1]%
  \restoreexpandmode
}

\newcommand{\splitatstar}[3]{%
  \protected@edef\split@temp{#1}%
  \saveexpandmode
  \expandarg\StrCut{\split@temp}{*}#2#3%
  \restoreexpandmode
}

\newcommand{\splitatinf}[3]{%
  \protected@edef\split@temp{#1}%
  \saveexpandmode
  \expandarg\StrCut{\split@temp}{<}#2#3%
  \restoreexpandmode
}

\newcommand{\splitatequal}[3]{%
  \protected@edef\split@temp{#1}%
  \saveexpandmode
  \expandarg\StrCut{\split@temp}{=}#2#3%
  \restoreexpandmode
}

\newcommand{\ifsubstring}[4]{%
  \protected@edef\split@temp{#1}%
  \protected@edef\split@tempbis{#2}%
  \saveexpandmode
  \expandarg\IfSubStr{\split@temp}{\split@tempbis}{#3}{#4}%
  \restoreexpandmode
}

\def\removespaces#1{\zap@space#1 \@empty}

%%% Commands for alphabetical ordering %%%

\newcommand{\sortitem}[2][\relax]{%
	\DTLnewrow{list}% Create a new entry
	\ifx#1\relax%
		\DTLnewdbentry{list}{sortlabel}{#2}% Add entry sortlabel (no optional argument)
	\else%
		\DTLnewdbentry{list}{sortlabel}{#1}% Add entry sortlabel (optional argument)
	\fi%
		\DTLnewdbentry{list}{description}{#2}% Add entry description
}
\newenvironment{sortedlist}{%
	\DTLifdbexists{list}{\DTLcleardb{list}}{\DTLnewdb{list}}% Create new/discard old list
}{%
	\DTLsort{sortlabel}{list}% Sort list
	\begin{itemize*}[label={}, itemjoin={,}]%
		\DTLforeach*{list}{\theDesc=description}{%
		\item\theDesc}% Print each item
	\end{itemize*}%
}

\pdfstringdefDisableCommands{\def\textcolor#1{}}

% See language specific file for \addtosortedlist

%%% Database for automatic Quick Ref Sheet %%%

\DTLnewdb{profiles} % Database containing name, category, multiprofile number, profilename (if multi), caraclist, trooptype, invocation for CV.
\newcommand{\profilecategory}{\labels@lords} % Will be updated in relevant categories

\newcommand{\profiledtbfillname}[1]{\DTLnewdbentry{profiles}{name}{#1}}
\newcommand{\profiledtbfillcategory}[1]{\DTLnewdbentry{profiles}{category}{#1}}
\newcommand{\profiledtbfilltrooptype}[1]{\DTLnewdbentry{profiles}{trooptype}{#1}}
\newcommand{\profiledtbfillinvocation}[1]{\DTLnewdbentry{profiles}{invocation}{#1}}
\newcommand{\profiledtbfillprofile}[1]{\DTLnewdbentry{profiles}{profile}{#1}}
\newcommand{\profiledtbfillmultipleprofile}[1]{\DTLnewdbentry{profiles}{multipleprofile}{#1}}

\newcommand{\void}[1]{}
\newcounter{multiprofilecounter}

\newcommand{\profiledtbfillcarac}[1]{%
	\profiledtbfillprofile{#1}
	\parselist{#1}{\locallists@profileslist}% Split of the different profiles in the case of a multiprofile.
	\setcounter{multiprofilecounter}{0}%
	\forlistloop{\stepcounter{multiprofilecounter}\void}{\locallists@profileslist}%
	\expandafter\profiledtbfillmultipleprofile\expandafter{\number\value{multiprofilecounter}}
}


%%% Technical commands %%%

\newcommand{\newrule}{\textcolor{green!50!black}}
\newcommand{\removedrule}[1]{\textcolor{green!50!black}{\sout{#1}}}
\newcommand{\starsymbol}{$\star$}
\newcommand{\refsymbol}{$^\star$}

\newcommand{\inch}{\arcsecond}
\newcommand{\foot}{\arcminute}
\newcommand{\range}[1] {\labels@range~\unit{#1}{\inch}}
\newcommand{\distance}[1] {\unit{#1}{\inch}}
\newcommand{\result}[1] {\texttt{'}#1\texttt{'}}
\newcommand{\plusone}{+1}


%%% Fonts and sizes %%%

\newcommand{\bigtitle}[1]{\vspace*{-1.5cm}\section*{}\noindent\begin{center}\Hugefontsize\textbf{\antiquefont\expandafter\uppercase\expandafter{#1}}\end{center}}

\newcommand{\subtitle}[1]{\subsection*{}\noindent{\hugefontsize\antiquefont #1}}

\newcommand{\subsubtitle}[1]{\subsubsection*{}\noindent{\Largerfontsize\antiquefont #1}}

\newcommand{\verysmallfontsize}{\fontsize{4}{4.8}\selectfont}
\newcommand{\smallfontsize}{\fontsize{6}{7.2}\selectfont}
\newcommand{\normalfontsize}{\fontsize{8}{9.6}\selectfont}
\newcommand{\largefontsize}{\fontsize{10}{12}\selectfont}
\newcommand{\largerfontsize}{\fontsize{12}{14.4}\selectfont}
\newcommand{\Largefontsize}{\fontsize{14}{16.8}\selectfont}
\newcommand{\Largerfontsize}{\fontsize{15}{18}\selectfont}
\newcommand{\hugefontsize}{\fontsize{18}{21.6}\selectfont}
\newcommand{\Hugefontsize}{\fontsize{25}{30}\selectfont}

\newcommand{\unitentryformat}[1]{\textit{\largefontsize{#1}}}
\newcommand{\textIT}[1]{\textit{\largefontsize{#1}}}


%%% Titles %%%

\newcommand{\lordstitle}{\def\logolocalpath{../Layout/pics/logo_lord.png}\bigtitle{\labels@lords}}
\newcommand{\heroestitle}{%
\def\logolocalpath{../Layout/pics/logo_hero.png}%
\clearpage\bigtitle{\labels@heroes}%
\renewcommand{\profilecategory}{\labels@heroes}%
}
\newcommand{\coreunitstitle}{%
\def\logolocalpath{../Layout/pics/logo_core.png}%
\clearpage\bigtitle{\labels@coreunits}%
\renewcommand{\profilecategory}{\labels@coreunits}%
}
\newcommand{\specialunitstitle}{%
\def\logolocalpath{../Layout/pics/logo_special.png}%
\clearpage\bigtitle{\labels@specialunits}%
\renewcommand{\profilecategory}{\labels@specialunits}%
}
\newcommand{\rareunitstitle}{%
\def\logolocalpath{../Layout/pics/logo_rare.png}%
\clearpage\bigtitle{\labels@rareunits}%
\renewcommand{\profilecategory}{\labels@rareunits}%
}
\newcommand{\mountstitle}{%
\def\logolocalpath{../Layout/pics/logo_mount.png}%
\clearpage\bigtitle{\labels@charactermounts}%
\renewcommand{\profilecategory}{\labels@mounts}%
}

\newcommand{\startarmywiderules}{\newpage\bigtitle{\labels@armywiderules}\largefontsize}
\newcommand{\closearmywiderules}{\normalfontsize}
\newcommand{\armywideruleentry}[1]{\subtitle{#1}\vspace{5pt}}

\newcommand{\startarmyspecialrules}{\bigtitle{\labels@armyspecialrules}\largefontsize}
\newcommand{\closearmyspecialrules}{\normalfontsize}
\newcommand{\armyspecialruleentry}[1]{\subtitle{#1}\vspace{5pt}}

\newcommand{\startarmyarmoury}{\bigtitle{\labels@armoury}\largefontsize\subtitle{}}
\newcommand{\closearmyarmoury}{\normalfontsize}

\newcommand{\startarmymagicalitems}{\newpage\largefontsize\bigtitle{\labels@magicalitems}\begin{multicols}{2}\raggedcolumns}
\newcommand{\closearmymagicalitems}{\end{multicols}\normalfontsize}

\newcommand{\armymagicalweapons}{\subtitle{\labels@magicalweapons}}
\newcommand{\armymagicalarmour}{\subtitle{\labels@magicalarmour}}
\newcommand{\armytalismans}{\subtitle{\labels@talismans}}
\newcommand{\armyenchanteditems}{\subtitle{\labels@enchanteditems}}
\newcommand{\armyarcaneitems}{\subtitle{\labels@arcaneitems}}
\newcommand{\armymagicalbanners}{\subtitle{\labels@magicalbanners}}

\newcommand{\startarmynewsection}[1]{\newpage\bigtitle{#1}\largefontsize}
\newcommand{\startarmynewsectionSP}[1]{\vspace{1.5cm}\bigtitle{#1}\largefontsize}
\newcommand{\closearmynewsection}{\normalfontsize}

\newcommand{\armynewsubsection}[1]{\subtitle{#1}\vspace{5pt}}
\newcommand{\armynewsubsubsection}[1]{\subsubtitle{#1}\vspace{3pt}}

\newcommand{\armylist}{\clearpage}

\newcommand{\quickrefsheettitle}{\clearpage\newgeometry{top=1.6cm, bottom=2cm, left=1cm, right=1cm}\bigtitle{\labels@quickrefsheet}\vspace*{0.4cm}}
\newcommand{\changelogtitle}{\clearpage\bigtitle{\labels@changelog}\spaceaftersection{}}

\newcommand{\spaceaftersection}{\vspace{0.8cm}}

\newcommand{\separator}{\noindent\begin{center}\textcolor{black!30}{\rule{0.7\columnwidth}{2pt}}\end{center}}


%%% Custom lists and description for first sections of the army books

\newcommand{\startpricelist}{\begin{samepage}\begin{description}[leftmargin=0.3cm, labelindent=0cm, labelsep=0.1cm]}
\def\endpricelist{\end{description}\end{samepage}}
\newcommand{\pricelistitem}[2]{\item \option{\textbf{#1}}{#2}\newline}
\newcommand{\nopricelistitem}[1]{\item \textbf{#1}\newline}

\newcommand{\startpricelistNSP}{\begin{description}[leftmargin=0.3cm, labelindent=0cm, labelsep=0.1cm]}
\def\endpricelistNSP{\end{description}}

\newcommand{\startitemlist}{\begin{multicols}{2}\raggedcolumns\begin{description}[leftmargin=0.3cm, labelindent=0cm, labelsep=0.1cm]}
\def\enditemlist{\end{description}\end{multicols}}
\newcommand{\listitem}[1]{\item[#1\spacebeforecolon{}:]}

\newcommand{\startitemlistonecol}{\begin{description}[leftmargin=0.3cm, labelindent=0cm, labelsep=0.1cm]}
\def\enditemlistonecol{\end{description}}
\newcommand{\listitemonecol}[1]{\item \textbf{#1\spacebeforecolon{}:}\newline}

\newenvironment{customitemize}{\begin{description}[leftmargin=0.3cm, labelindent=0cm, labelsep=0cm]}{\end{description}}
\newenvironment{customsubitemize}{\begin{itemize}[label={-}, labelsep=0.1cm, topsep=0cm, parsep=0cm, itemsep=0cm, leftmargin=0.4cm, labelindent=0cm]}{\end{itemize}}

%%% Table parameters %%%

\newcolumntype{M}[1]{>{\centering\let\newline\\\arraybackslash\hspace{0pt}}m{#1}}


%%%  Lists handling %%%

\newcommand{\addlocallist}{\listadd\locallists@dummy}%
\NewDocumentCommand{\parsespacelist}{>{\SplitList{ }} m }{%
	\ProcessList{#1}{\addlocallist}%
}%
\NewDocumentCommand{\parsecommalist}{>{\SplitList{,}} m }{%
	\ProcessList{#1}{\addlocallist}%
}%
\newcommand{\parselist}[3][,]{%
	\renewcommand\addlocallist{\listadd#3}%
  	\undef#3%
  	\ifstrequal{#1}{ }{\parsespacelist{#2}}{\parsecommalist{#2}}%
}


%%% Profiles handling %%%

% Element of a table that contains the characteristics of a model (or part of a model)
\newcommand\caraclist[1]{
	\parselist[ ]{#1}{\locallists@caraclist}%
	\forlistloop{&}{\locallists@caraclist}%
}

\newcommand\caraclistbold[1]{
	\parselist[ ]{#1}{\locallists@caraclist}%
	\forlistloop{&\bfseries}{\locallists@caraclist}%
}

% Line of a profile table, including bottom line. It is meant to contain the name of the model (or part), its characteristics (preferably, the second argument should contain the \carac macro), troop type and base size.
\newcommand{\profilefirstline}[4]{#1 & #2 &   & #3 & #4 }

% Start of a profile table. Includes the table commands, and the column labels. \profilecellsize is the size of the characteristics cells in the profile.
\newcommand{\profilecellsize}{0.56cm}
\newcommand{\profilestart}{%
	\noindent %
	\begin{tabular}{@{}p{3cm}@{}M{\profilecellsize}@{}M{\profilecellsize}@{}M{\profilecellsize}@{}M{\profilecellsize}@{}M{\profilecellsize}@{}M{\profilecellsize}@{}M{\profilecellsize}@{}M{\profilecellsize}@{}M{\profilecellsize}@{}p{2.7cm}@{}p{3.3cm}@{}p{2cm}@{}}%
	 &% \textbf{\labels@profile}
	\labels@M & \labels@WS & \labels@BS & \labels@S & \labels@T & \labels@W & \labels@I & \labels@A & \labels@Ld &%
	&%
	{\unitentryformat{\labels@trooptype}} &%
	{\unitentryformat{\labels@basesize}}%
}

% End of a profile table.
\newcommand{\profileend}{\end{tabular}}

% Algorithm to automatically use and fill previous command, with coherence check.
\providebool{profilefirst}
\newcommand{\profileitem}[1]{%
	\tabularnewline%
	\splitatinf{#1}\local@unitname\local@unitprofile%
	\local@unitname \expandafter\caraclistbold\expandafter{\local@unitprofile}%
	&%
	& \ifbool{profilefirst}{\unit@type}{}%
	& \ifbool{profilefirst}{%
		\ifsubstring{\unit@basesize}{x}{% Rectangular base
			\unit{\unit@basesize}{\milli\meter}%
		}{% Circular base
			\unit{\unit@basesize}{\milli\meter} \labels@roundbase%
		}%
	}{}%
	\global\boolfalse{profilefirst}%
}
\newcommand{\profile}[1]{%
	\parselist{#1}{\locallists@profileslist}%
	\profilestart%
	\global\booltrue{profilefirst}%
	\forlistloop{\profileitem}{\locallists@profileslist}%
	\profileend%
}


%%% Profiles handling in case of invocation %%%

\newcommand{\invocprofilestart}{%
	\noindent %
	\begin{tabular}{@{}p{3cm}@{}M{\profilecellsize}@{}M{\profilecellsize}@{}M{\profilecellsize}@{}M{\profilecellsize}@{}M{\profilecellsize}@{}M{\profilecellsize}@{}M{\profilecellsize}@{}M{\profilecellsize}@{}M{\profilecellsize}@{}M{2.2cm}@{}p{0.5cm}@{}p{3.3cm}@{}p{2cm}@{}}%
	 &% \textbf{\labels@profile}
	\labels@M & \labels@WS & \labels@BS & \labels@S & \labels@T & \labels@W & \labels@I & \labels@A & \labels@Ld & \unitentryformat{\labels@Invocation} &%
	&%
	{\unitentryformat{\labels@trooptype}} &%
	{\unitentryformat{\labels@basesize}}%
}

\newcommand{\invocprofileitem}[1]{%
	\tabularnewline%
	\splitatinf{#1}\local@unitname\local@unitprofile%
	\local@unitname \expandafter\caraclistbold\expandafter{\local@unitprofile}%
	& \ifbool{profilefirst}{\unit@invocation}{} &%
	& \ifbool{profilefirst}{\unit@type}{}%
	& \ifbool{profilefirst}{\unit{\unit@basesize}{\milli\meter}}{}%
	\global\boolfalse{profilefirst}%
}

\newcommand{\invocprofile}[1]{%
	\parselist{#1}{\locallists@profileslist}%
	\invocprofilestart%
	\global\booltrue{profilefirst}%
	\forlistloop{\invocprofileitem}{\locallists@profileslist}%
	\profileend%
}


%%%%%%%%%%%%%%%%%%
%%% Unit rules %%%
%%%%%%%%%%%%%%%%%%

%%% Entry title command %%%

\newcommand{\unitentry}[2]{\ifdefempty{#1}{}{\noindent #2}}


%%% Special rules %%%

% Special rules listing for a unit, with alphabetical order.
\newcommand{\ruleslist}[1]{%
	\parselist[,]{#1}{\locallists@ruleslist}%
	\begin{sortedlist}%
		\forlistloop{\addtosortedlist}{\locallists@ruleslist}%
	\end{sortedlist}%
}

% Special rules entry.
\newcommand{\specialrules}[1]{\unitentry{#1}{\unitentryformat{\labels@specialrules\spacebeforecolon{}:}\newline\hspace*{-\fontdimen2\font}\expandafter\ruleslist\expandafter{#1}.}}
\newcommand{\commonspecialrules}[2]{\unitentry{#2}{\unitentryformat{#1\spacebeforecolon{}:}\newline\hspace*{-\fontdimen2\font}\expandafter\ruleslist\expandafter{#2}.}}


%%% Magical abilities %%%

% Paths listing for a unit.
\newcommand{\pathslist}[1]{%
	\parselist[,]{#1}{\locallists@pathslist}%
	\begin{itemize*}[label={}, itemjoin={,}, itemjoin*={\listlastchoice}]%
		\forlistloop{\item}{\locallists@pathslist}%
	\end{itemize*}%
}

% Magic entry.
\newcommand{\magic}[2]{\unitentry{#2}{\unitentryformat{\labels@magic\spacebeforecolon{}: }\newline\ifdefempty{#1}{}{\textbf{\magiclevel{#1}}. }\labels@pathsused\expandafter\pathslist\expandafter{#2}.}}

% Wizard Conclave.
\newcommand{\magicwizardconclave}[1]{\unitentry{#1}{\unitentryformat{\labels@magic\spacebeforecolon{}: }\newline\textbf{\wizardconclave{}}\spacebeforecolon{}: #1.}}


%%% Equipment %%%

% Equipment listing.
\newcommand{\equipmentlist}[1]{%
	\parselist[,]{#1}{\locallists@equipmentlist}%
	\begin{sortedlist}%
		\forlistloop{\addtosortedlist}{\locallists@equipmentlist}%
	\end{sortedlist}%
}

% Equipment entry.
\newcommand{\weapons}[1]{\unitentry{#1}{\unitentryformat{\labels@weapons\spacebeforecolon{}:}\newline\hspace*{-\fontdimen2\font}\expandafter\equipmentlist\expandafter{#1}.}}

\newcommand{\armour}[1]{\unitentry{#1}{\unitentryformat{\labels@armour\spacebeforecolon{}:}\newline\hspace*{-\fontdimen2\font}\expandafter\equipmentlist\expandafter{#1}.}}


%%% Alignment %%%

\newcommand{\alignment}[1]{\unitentry{#1}{\unitentryformat{\labels@alignment\spacebeforecolon{}:}\newline\textbf{#1}.}}

%%% Green Hide Race %%%

\newcommand{\greenhideraceentry}[1]{\unitentry{#1}{\unitentryformat{\labels@greenhiderace\spacebeforecolon{}:}\newline\textbf{#1}.}}


%%% Options %%%

% Frame commands.
\newcommand{\optionsframestart}{\begin{innerframe}[\labels@options]}
\newcommand{\optionsframeend}{\end{innerframe}}

% Options listing.
\newcommand{\optionslist}[1]{%
	\parselist[,]{#1}{\locallists@optionslist}%
	\begin{description}[leftmargin=0.3cm, labelindent=0cm, labelsep=0cm, itemsep=0cm, parsep=0cm]%
		\forlistloop{\item\setoption}{\locallists@optionslist}%
	\end{description}%
}

% Options entry.
\newcommand{\options}[1]{\ifdefempty{#1}{}{\optionsframestart\vspace*{-0.4cm}\unitentry{#1}{\expandafter\optionslist\expandafter{#1}}\optionsframeend}}

% Option specific commands.
\newcommand{\setoption}[1]{%
	\noexpandarg\StrCut{#1}{=}\optiontext\optionvalue%
	\expandafter\ifstrequal\expandafter{\optionvalue}{}{%
		\optiontext%
	}{%
	\ifsubstring{\optionvalue}{\free}{%
		\option[\free]{\optiontext}{\optionvalue}%
	}{%
	\ifsubstring{\optionvalue}{\unlimited}{%
		\option[\unlimited]{\optiontext}{\optionvalue}%
	}{%
	\ifsubstring{\optionvalue}{\upto}{%
		\splitatinf{\optionvalue}\myoption\myvalue%
		\option[\upto]{\optiontext}{\myvalue}%
	}{%
	\ifsubstring{\optionvalue}{\permodel}{%
		\splitatinf{\optionvalue}\myoption\myvalue%
		\option[\permodel]{\optiontext}{\myvalue}%
	}{%
	\ifsubstring{\optionvalue}{\pershadygit}{% For Orcs N Goblins
		\splitatinf{\optionvalue}\myoption\myvalue%
		\option[\pershadygit]{\optiontext}{\myvalue}%
	}{%
	\ifsubstring{\optionvalue}{\permadgit}{% For Orcs N Goblins
		\splitatinf{\optionvalue}\myoption\myvalue%
		\option[\permadgit]{\optiontext}{\myvalue}%
	}{%	
	\ifsubstring{\optionvalue}{\perrune}{% For Dwarven Holds
		\splitatinf{\optionvalue}\myoption\myvalue%
		\option[\perrune]{\optiontext}{\myvalue}%
	}{%	
		\option{\optiontext}{\optionvalue}%
	}}}}}}}}%
}

\newcommand{\option}[3][]{#2\predotfill\dotfill\nobreak%
	% Add \upto token if necessary.
	\ifstrequal{#1}{\upto}{\upto~}{}%
	% The option can be free, have an unlimited cost, or have a points cost.
	\ifstrequal{#1}{\free}{\free}{\ifstrequal{#1}{\unlimited}{\unlimited}{\pts{#3}}}%
	% Add \permodel if necessary.
	\ifstrequal{#1}{\permodel}{\nobreak\permodel}{}%
	% Add \persomething if necessary.
	\ifstrequal{#1}{\pershadygit}{\nobreak\pershadygit}{}% For Orcs N Goblins
	\ifstrequal{#1}{\permadgit}{\nobreak\permadgit}{}% For Orcs N Goblins
	\ifstrequal{#1}{\perrune}{\nobreak\perrune}{}% For Dwarven Holds
}

\newcommand\optionschoice[2]{%
	\parselist[,]{#2}{\locallists@optionschoice}%
	#1%
	\begin{itemize}[label={}, parsep=0cm, labelindent=0cm, labelwidth=0cm, noitemsep, topsep=0em, leftmargin=0.3cm]%
	\forlistloop{\item\setoption}{\locallists@optionschoice}%
	\end{itemize}%
}

\newcommand\optionschoiceTWOCOL[2]{%
	\parselist[,]{#2}{\locallists@optionschoice}%
	#1%
	\begin{itemize}[label={}, parsep=0cm, labelindent=0cm, labelwidth=0cm, noitemsep, topsep=0em, leftmargin=0.3cm]%
	\setlength{\columnseprule}{0.5pt}
	\renewcommand{\columnseprulecolor}{\color{black!30}}
	\vspace*{-5pt}\begin{multicols}{2}\raggedcolumns
	\forlistloop{\item\setoption}{\locallists@optionschoice}%
	\end{multicols}\setlength{\columnseprule}{0pt}
	\end{itemize}%
}

% Option description in army desc.
\newcommand{\optiondef}[3]{\option{\textbf{#1}}{#2}\ifblank{#3}{}{\\{#3}}}


%%% Mount options %%%

% Frame commands.
\newcommand{\mountsframestart}{\begin{innerframe}[\labels@mounts]}
\newcommand{\mountsframeend}{\end{innerframe}}

% Mount listing.
\newcommand{\mountslist}[1]{%
	\parselist[,]{#1}{\locallists@mountslist}%
	\begin{description}[leftmargin=0.3cm, labelindent=0cm, labelsep=0cm, itemsep=0cm, parsep=0cm]%
		\forlistloop{\item\setoption}{\locallists@mountslist}%
	\end{description}%
}

% Mount entry.
\newcommand{\mounts}[1]{\ifdefempty{#1}{}{\mountsframestart\vspace*{-0.4cm}\unitentry{#1}{\expandafter\mountslist\expandafter{#1}}\mountsframeend}}


%%% Command group %%%

% Command group specific commands.
\define@key{commandgroup}{restriction}            {\def\commandgroup@restriction{#1}}
\define@key{commandgroup}{champion}               {\def\commandgroup@champion{#1}}
\define@key{commandgroup}{championallowance}      {\def\commandgroup@championallowance{#1}}
\define@key{commandgroup}{championoption}         {\def\commandgroup@championoption{#1}}
\define@key{commandgroup}{championprerestriction} {\def\commandgroup@championprerestriction{#1}}
\define@key{commandgroup}{championrestriction}    {\def\commandgroup@championrestriction{#1}}
\define@key{commandgroup}{banner}                 {\def\commandgroup@banner{#1}}
\define@key{commandgroup}{bannerallowance}        {\def\commandgroup@bannerallowance{#1}}
\define@key{commandgroup}{veteranstandardbearer}  {\def\commandgroup@veteranstandardbearer{#1}}
\define@key{commandgroup}{singlebannerallowance}  {\def\commandgroup@singlebannerallowance{#1}}
\define@key{commandgroup}{condsinglebannerallowance}  {\def\commandgroup@condsinglebannerallowance{#1}}
\define@key{commandgroup}{banneroption}           {\def\commandgroup@banneroption{#1}}
\define@key{commandgroup}{bannerrestriction}      {\def\commandgroup@bannerrestriction{#1}}
\define@key{commandgroup}{musician}               {\def\commandgroup@musician{#1}}
\define@key{commandgroup}{musicianrestriction}    {\def\commandgroup@musicianrestriction{#1}}
\newcommand{\defcommandgroup}{%
	\setkeys{commandgroup}{restriction=,
	                       champion=, championallowance=, championoption=, championprerestriction=, 
	                       championrestriction=, banner=, bannerallowance=, veteranstandardbearer=, 
	                       singlebannerallowance=, condsinglebannerallowance=, banneroption=, 
	                       bannerrestriction=, musician=, musicianrestriction=}%
	\setkeys{commandgroup}%
}

% Frame commands.
\newcommand{\commandgroupframestart}{\begin{innerframe}[\labels@commandgroup]}
\newcommand{\commandgroupframeend}{\end{innerframe}}

% Command group entry.
\newcommand{\commandgroup}[1]{%
	\defcommandgroup{#1}%
	\ifstrempty{#1}{}{\commandgroupframestart\vspace*{-0.2cm}%
		\begin{description}[leftmargin=0.3cm, labelindent=0cm, labelsep=0cm, itemsep=0cm, parsep=0cm]%
			% Command group title, including restrictions applying to all the command group
			\item \textbf{\expandafter\ifblank\expandafter{\commandgroup@restriction}{}{ \only{\commandgroup@restriction}\spacebeforecolon{}: }} 
			% Champion handling.
			\ifdefempty{\commandgroup@champion}{}{% We have a champion!
			\ifdefempty{\commandgroup@championprerestriction}{% There is no prerestriction to have a champion
				\item \hspace*{-0.04cm}\option{\labels@champion%
					% Possible restrictions to taking a champion
				    \expandafter\ifblank\expandafter{\commandgroup@championrestriction}{}{ \only{\commandgroup@championrestriction}}%
				    % Cost of a champion
				    }{\commandgroup@champion}%
				    % Magical allowance of the champion. Should probably not be used, champion option can do it as well and is more flexible.
					\ifdefempty{\commandgroup@championallowance}{}{\par\option[\upto]{\hspace*{0.3cm}- \labels@championallowance}{\commandgroup@championallowance}}%
					% Any option available to the champion, in the form option:cost
					\ifdefempty{\commandgroup@championoption}{}{%
						\splitatinf{\commandgroup@championoption}\local@option\local@cost%
						\par\option{\hspace*{0.3cm}- \local@option}{\local@cost}}%
			}{% There is a pre-restriction to have a champion
				\item \hspace*{-0.04cm}\commandgroup@championprerestriction	\newline%
				\option{\labels@champion}{\commandgroup@champion}%
				% Magical allowance of the champion. Should probably not be used, champion option can do it as well and is more flexible.
				\ifdefempty{\commandgroup@championallowance}{}{\par\option[\upto]{\hspace*{0.3cm}- \labels@championallowance}{\commandgroup@championallowance}}%
				% Any option available to the champion, in the form option:cost
				\ifdefempty{\commandgroup@championoption}{}{%
					\splitatinf{\commandgroup@championoption}\local@option\local@cost%
					\par\option{\hspace*{0.3cm}- \local@option}{\local@cost}}%
			} %End of the prerestriction of not condition
			}% End of champion handling
			\ifdefempty{\commandgroup@musician}{}{% We have a musician!
				\item \hspace*{-0.04cm}\option{\labels@musician%
					% Possible restrictions to taking a musician
				    \expandafter\ifblank\expandafter{\commandgroup@musicianrestriction}{}{ \only{\commandgroup@musicianrestriction}}%
				    % Cost of a musician
				    }{\commandgroup@musician}%
			}%
			\ifdefempty{\commandgroup@banner}{}{% We have a banner!
				\item \hspace*{-0.04cm}\option{\labels@standardbearer%
					% Possible restrictions to taking a banner
				    \expandafter\ifblank\expandafter{\commandgroup@bannerrestriction}{}{ \only{\commandgroup@bannerrestriction}}%
				    % Cost of a banner
				    }{\commandgroup@banner}%
				    % Magical banner, if all units of this type can take one.
					\ifdefempty{\commandgroup@bannerallowance}{}{\par\option[\upto]{\hspace*{0.3cm}- \labels@bannerallowance}{\commandgroup@bannerallowance}}%
					% Magical banner, if Veteran.
					\ifdefempty{\commandgroup@veteranstandardbearer}{}{\par\hspace*{0.3cm}- \labels@veteranstandardbearer%
					\expandafter\ifstrequal\expandafter{\commandgroup@veteranstandardbearer}{*}{*}{}%
					}%
					% Magical banner, if only one unit of this type can take one.
					\ifdefempty{\commandgroup@singlebannerallowance}{}{\par\option[\upto]{\hspace*{0.3cm}- \labels@singlebannerallowance}{\commandgroup@singlebannerallowance}}%
					% Magical banner, if only one unit of this type can take one, but with condtions.
					\ifdefempty{\commandgroup@condsinglebannerallowance}{}{%
						\splitatinf{\commandgroup@condsinglebannerallowance}\local@option\local@cost%
						\par\option[\upto]{\hspace*{0.3cm}- \labels@condsinglebannerallowance \local@option}{\local@cost}}%
					% Additional option for the banner, such as Hill Goblin Lookouts for Ogres
					\ifdefempty{\commandgroup@banneroption}{}{%
						\splitatinf{\commandgroup@banneroption}{\local@option}{\local@cost}%
						\par\option{\hspace*{0.3cm}- \local@option}{\local@cost}%
					}%
			}%
		\end{description}%
	\commandgroupframeend%
	 }%
}


%%% Unit rules %%%

% Frame commands.
\newcommand{\unitrulesframestart}{\begin{innerframe}[\labels@specialrules]}
\newcommand{\unitrulesframeend}{\end{innerframe}}

% Unit rules specific commands.
\newcommand{\unitrule}[2]{\item[#1\spacebeforecolon{}:]#2}

% Unit rule entry.
\newcommand{\unitrules}[1]{\ifdefempty{#1}{}{\unitrulesframestart\vspace*{-0.05cm}\begin{description}[leftmargin=0.3cm, labelindent=0cm, labelsep=0.1cm, itemsep=0.2cm, parsep=0cm]#1\end{description}\unitrulesframeend}}


%%% Special equipment %%%

% Frame commands.
\newcommand{\unitequipmentframestart}{\begin{innerframe}[\labels@specialequipment]}
\newcommand{\unitequipmentframeend}{\end{innerframe}}

% Special equipment specific commands.
\newcommand{\equipmentdef}[2]{\item[#1\spacebeforecolon{}:]#2}

% Special equipment entry.
\newcommand{\unitequipment}[1]{\ifdefempty{#1}{}{\unitequipmentframestart\vspace*{-0.05cm}\begin{description}[leftmargin=0.3cm, labelindent=0cm, labelsep=0.1cm, itemsep=0.2cm, parsep=0cm]#1\end{description}\unitequipmentframeend}}






%%%%%%%%%%%%%%%%%%%%%%%%%%%%%%%%
%%% Profile input and layout %%%
%%%%%%%%%%%%%%%%%%%%%%%%%%%%%%%%

%%% Input parameters %%%

\define@key{unit}{notinQRS}{\def\unit@notinQRS{#1}}
\define@key{unit}{name}{\def\unit@name{#1}}
\define@key{unit}{QRSname}{\def\unit@QRSname{#1}}
\define@key{unit}{profile}{\def\unit@profile{#1}}
\define@key{unit}{cost}{\def\unit@cost{#1}}
\define@key{unit}{invocation}{\def\unit@invocation{#1}}
\define@key{unit}{costpermodel}{\def\unit@costpermodel{#1}}
\define@key{unit}{maxmodels}{\def\unit@maxmodels{#1}}
\define@key{unit}{type}{\def\unit@type{#1}}
\define@key{unit}{unitsize}{\def\unit@unitsize{#1}}
\define@key{unit}{basesize}{\def\unit@basesize{#1}}
\define@key{unit}{commonspecialrules}{\def\unit@commonspecialrules{#1}}
\define@key{unit}{commontype}{\def\unit@commontype{#1}}
\define@key{unit}{commonspecialrulesB}{\def\unit@commonspecialrulesB{#1}}
\define@key{unit}{commontypeB}{\def\unit@commontypeB{#1}}
\define@key{unit}{specialrules}{\def\unit@specialrules{#1}}
\define@key{unit}{magiclevel}{\def\unit@magiclevel{#1}}
\define@key{unit}{magicpaths}{\def\unit@magicpaths{#1}}
\define@key{unit}{equipment}{\def\unit@equipment{#1}}
\define@key{unit}{alignment}{\def\unit@alignment{#1}}
\define@key{unit}{greenhiderace}{\def\unit@greenhiderace{#1}}
\define@key{unit}{weapons}{\def\unit@weapons{#1}}
\define@key{unit}{armour}{\def\unit@armour{#1}}
\define@key{unit}{wizardconclave}{\def\unit@wizardconclave{#1}}
\define@key{unit}{unitequipment}{\def\unit@unitequipment{#1}}
\define@key{unit}{options}{\def\unit@options{#1}}
\define@key{unit}{mounts}{\def\unit@mounts{#1}}
\define@key{unit}{commandgroup}{\def\unit@commandgroup{#1}}
\define@key{unit}{unitrules}{\def\unit@unitrules{#1}}
\define@key{unit}{additional}{\def\unit@additional{#1}}


%%% Frames definition %%%

% Unit's big frame.
\tikzset{unitprice/.style={draw=white, fill=white, rectangle, rounded corners, right, minimum height=0.7cm}}
\tikzset{unittitle/.style={draw=white, fill=white, rectangle, rounded corners, right, minimum height=0.7cm, font=\bfseries}}
\tikzset{unitlogo/.style={draw=white, fill=white, rectangle, right, minimum height=0.7cm}}

\newenvironment{unitframe}[2][]{%
	\mdfsetup{%
		nobreak=true,%
		linewidth=1pt,%
		linecolor=black!30,%
		roundcorner=5pt,%
		backgroundcolor=white,%
		innertopmargin=1.2\baselineskip,
		innerbottommargin=1.2\baselineskip,
		singleextra={
			\expandafter\ifblank\expandafter{\unit@cost}{}{%
				\node[unitprice,anchor=east,xshift=-0.5cm] at (P)%
					{%
						{{\smallfontsize\minprice} \Largefontsize\pts{\textbf{\unit@cost}}}%
					};
				}%
				\node[unittitle,xshift=0.5cm] at (P-|O)%
					{\Largefontsize\antiquefont\uppercase\expandafter\expandafter\expandafter{\unit@name}};
				\node[unitlogo, xshift=8.1cm, yshift=0.1cm] at (P-|O)%
					{\includegraphics[width=1.2cm]{\logolocalpath}};
		}
	}%
	\begin{mdframed}[]\relax%
}%
{%
\end{mdframed}%
}

% Inner small frames for options, special rules definition, ...
\tikzset{innertitle/.style={fill=white, rectangle, rounded corners, right, minimum height=8pt, xshift=0.5cm}}

\newenvironment{innerframe}[1][]{%
	\mdfsetup{%
		innerleftmargin=5pt,%
		innerrightmargin=5pt,%
		linecolor=black!30,%
		linewidth=0.5pt,%
		roundcorner=5pt,%
		backgroundcolor=white,%
		innertopmargin=1.1\baselineskip,
		singleextra={
		\node[innertitle] at (P-|O)%
			{\unitentryformat{#1}};
		}
	}%
	\vspace*{-0.2cm}\begin{mdframed}[]\relax%
}%
{%
\end{mdframed}%
}

%%% Command to add a new unit definition %%%

\newcommand{\defunit}{
	\setkeys{unit}{%
		notinQRS=, name=, QRSname=, profile=, cost=, invocation=, costpermodel=, maxmodels=, type=, unitsize=, basesize=, commonspecialrules=, commontype=, commonspecialrulesB=, commontypeB=, specialrules=, magiclevel=, magicpaths=, alignment=, greenhiderace=, equipment=, weapons=, armour=, wizardconclave=, unitequipment=, options=, mounts=, commandgroup=, unitrules=, additional=%
	}%
	\setkeys{unit}%
}

\newcommand{\showunit}[1]{
	\defunit{#1}
	\begin{unitframe}[\unit@name]{\unit@cost}
	\mdfsetup{style=defaultoptions}
	\expandafter\ifblank\expandafter{\unit@unitsize}{}{%
	\expandafter\ifstrequal\expandafter{\unit@unitsize}{1}{% single model
		% Can you add model to this single model ?
		\expandafter\ifblank\expandafter{\unit@maxmodels}{% no		
			{\hspace*{0.25cm}\labels@Singlemodel}%
		}{% yes
			{\hspace*{0.25cm}\mincostfor{} \textbf{1} \labels@model{}. \maxunitsize{}\spacebeforecolon{}: \textbf{\unit@maxmodels} \labels@models{}.\hfill \additionalfigscost{} {\largefontsize\pts{\textbf{\unit@costpermodel{}}}\permodel}\hspace*{0.1cm}}%
		}%
	}{% not single model
		% Test if we wanna print a sentence instead of unit number
		\ifsubstring{\unit@unitsize}{SPECIAL-}{%
			\hspace*{0.25cm}\StrDel{\unit@unitsize}{SPECIAL-}%
		}{%	
			{\hspace*{0.25cm}\mincostfor{} \textbf{\unit@unitsize} \labels@models{}. \maxunitsize{}\spacebeforecolon{}: \textbf{\unit@maxmodels} \labels@models{}.\hfill \additionalfigscost{} {\largefontsize\pts{\textbf{\unit@costpermodel{}}}\permodel}\hspace*{0.1cm}}%
		}%
	}%
	}%
	\vspace*{-0.1cm}
	\noindent\begin{center}\textcolor{black!30}{\rule{\columnwidth}{1pt}}\end{center}
		\expandafter\ifblank\expandafter{\unit@invocation}{%
			\expandafter\profile\expandafter{\unit@profile}
		}{%
			\expandafter\invocprofile\expandafter{\unit@profile}
		}
	\noindent\begin{center}\textcolor{black!30}{\rule{\columnwidth}{1pt}}\end{center}
	\vspace*{-0.2cm}
	\setlength\multicolsep{0pt}
	\begin{multicols}{2}
		\raggedcolumns
		\vspace*{-0.3cm}{\setlength{\parskip}{0.3cm}
		\expandafter\ifblank\expandafter{\unit@alignment}{}{\noindent\parbox{\columnwidth}{\alignment{\unit@alignment}}}
		
		\expandafter\ifblank\expandafter{\unit@greenhiderace}{}{\noindent\parbox{\columnwidth}{\greenhideraceentry{\unit@greenhiderace}}}
		
		\expandafter\ifblank\expandafter{\unit@equipment}{}{\noindent\parbox{\columnwidth}{\equipment{\unit@equipment}}}
				
		\expandafter\ifblank\expandafter{\unit@weapons}{}{\noindent\parbox{\columnwidth}{\weapons{\unit@weapons}}}
		
		\expandafter\ifblank\expandafter{\unit@armour}{}{\noindent\parbox{\columnwidth}{\armour{\unit@armour}}}
		
		\expandafter\ifblank\expandafter{\unit@commonspecialrules}{}{\noindent\parbox{\columnwidth}{\commonspecialrules{\unit@commontype}{\unit@commonspecialrules}}}
		
		\expandafter\ifblank\expandafter{\unit@commonspecialrulesB}{}{\noindent\parbox{\columnwidth}{\commonspecialrules{\unit@commontypeB}{\unit@commonspecialrulesB}}}
		
		\expandafter\ifblank\expandafter{\unit@specialrules}{}{\noindent\parbox{\columnwidth}{\specialrules{\unit@specialrules}}}
		
		\expandafter\ifblank\expandafter{\unit@magicpaths}{}{\noindent\parbox{\columnwidth}{\magic{\unit@magiclevel}{\unit@magicpaths}}}
		
		\expandafter\ifblank\expandafter{\unit@wizardconclave}{}{\noindent\parbox{\columnwidth}{\magicwizardconclave{\unit@wizardconclave}}}
		}
		\vspace{0.1cm}
		\mounts{\unit@mounts}
		\options{\unit@options}
		\expandafter\ifblank\expandafter{\unit@commandgroup}{}{\expandafter\commandgroup\expandafter{\unit@commandgroup}}
		\unitrules{\unit@unitrules}
		\unitequipment{\unit@unitequipment}
	\end{multicols}
	\vspace*{0.1cm}\unit@additional
	\end{unitframe}
	% Database filling for auto QRS
	\expandafter\ifblank\expandafter{\unit@notinQRS}{%
	\DTLnewrow{profiles}%
	\expandafter\ifblank\expandafter{\unit@QRSname}{%
		\expandafter\profiledtbfillname\expandafter{\unit@name}%
	}{%
		\expandafter\profiledtbfillname\expandafter{\unit@QRSname}%
	}
	\expandafter\profiledtbfillcategory\expandafter{\profilecategory}%
	\expandafter\profiledtbfilltrooptype\expandafter{\unit@type}%
	\expandafter\ifblank\expandafter{\unit@invocation}{}{\expandafter\profiledtbfillinvocation\expandafter{\unit@invocation}}%
	\expandafter\profiledtbfillcarac\expandafter{\unit@profile}
	}{}%
}


%%% Changelog commands %%%

\newcommand{\newlog}[2]{%
\vspace*{0.2cm}\noindent{\antiquefont\Large\textbf{V#1}}
\parselist[,]{#2}{\locallists@changelist}%
\begin{itemize}[itemsep=0pt]%
\forlistloop{\item[-]}{\locallists@changelist}%
\end{itemize}%
}

\newcommand{\startchangelog}{\begin{multicols}{2}\vspace*{-0.2cm}}
\def\endchangelog{\end{multicols}}


\newcommand{\booktitle}{Forteresses Naines}
\newcommand{\version}{0.99.9}
\newcommand{\frenchversion}{2.1}
\newcommand{\booklogo}{pics/logo_big_DH.png}
\newcommand{\tocfirstcolumn}{%
\tocentry{widerulestitle}{\labels@armywiderules}

\tocentry{specialrulestitle}{\labels@armyspecialrules}

\tocentry{armourytitle}{\labels@armoury}

\tocentry{battlerunestitle}{\battlerunes}

\tocentry{dwarvenrunestitle}{\dwarvenrunes}

\tocentry{lordtitle}{\labels@armylist}

\tocentry{QRStitle}{\labels@quickrefsheet}
}
\hypersetup{pdftitle={T9A - \booktitle}}

% Army wide rules

\newcommand{\hewnoutofmountains}{Taillés dans les Montagnes}
\newcommand{\ancientgrudge}{Vieille Rancune}
\newcommand{\ancientgrudges}{Vieilles Rancunes}

% Army special rules

\newcommand{\relentless}{Implacable}
\newcommand{\sturdy}{Robuste}
\newcommand{\shieldwall}{Mur de Boucliers}
\newcommand{\entrench}{Retranchement}
\newcommand{\holdstone}{Pierre de la Forteresse}
\newcommand{\thebiggertheyare}{Plus Ils Sont Gros\dots{}}
\newcommand{\comewithme}{Tu Partiras avec Moi !}

% Army common type special rules

\newcommand{\dwarvencommonrules}{Règles Spéciales des Nains}
\newcommand{\seekerscommonrules}{Règles Spéciales des Chasseurs}

% Armoury

\newcommand{\wyrmslayerrocket}{Roquettes Tueuses de Wyrm}
\newcommand{\forgerepeater}{Sulfateuse}

% Battle Runes

\newcommand{\battlerunes}{Runes de Bataille}

\newcommand{\runeofmetal}{Rune de Métal}
\newcommand{\runeofresolve}{Rune de Résolution}
\newcommand{\runeofreckoning}{Rune de Jugement}
\newcommand{\runeofoaths}{Rune des Serments}
\newcommand{\runeofgleaming}{Rune Luisante}
\newcommand{\runeofresilience}{Rune de Résilience}

% Dwarven Runes

\newcommand{\dwarvenrunes}{Runes Naines}

\newcommand{\runeofpenetrating}{Rune de Pénétration}
\newcommand{\runeofdestruction}{Rune de Destruction}
\newcommand{\runeofsmashing}{Rune Brisante}
\newcommand{\runeofcraftsmanship}{Rune de l'Art de la Forge}
\newcommand{\runeofmight}{Rune de Puissance}
\newcommand{\runeofprecision}{Rune de Précision}
\newcommand{\runeofquickening}{Rune d'Accélération}
\newcommand{\runeoffury}{Rune de Fureur}
\newcommand{\runeoffire}{Rune de Feu}

\newcommand{\runeofresistance}{Rune de Résistance}
\newcommand{\runeofbronze}{Rune de Bronze}
\newcommand{\runeofsteel}{Rune d'Acier}
\newcommand{\runeofmountains}{Rune des Montagnes}
\newcommand{\runeofaegis}{Rune de l'Égide}
\newcommand{\runeofiron}{Rune de Fer}
\newcommand{\runeofcrushing}{Rune d'Écrasement}

\newcommand{\runeofretribution}{Rune de Châtiment}
\newcommand{\runeofdragonsbreath}{Rune du Dragon}
\newcommand{\runeofcourage}{Rune de Courage}
\newcommand{\runeofinfamy}{Rune d'Infamie}
\newcommand{\runeofshielding}{Rune Bouclier}
\newcommand{\runeofkinship}{Rune de Fraternité}
\newcommand{\runeoftheforge}{Rune de la Forge}

\newcommand{\runeofgrounding}{Rune de Connaissance}
\newcommand{\runeofmastery}{Rune de Maîtrise}
\newcommand{\runeofdevouring}{Rune Dévoreuse}
\newcommand{\runeofdenial}{Rune de Déni}
\newcommand{\runeofbrilliance}{Rune Brillante}
\newcommand{\runeofharnessing}{Rune d'Exploitation}

\newcommand{\runicstandardofshielding}{Bannière Runique de Protection}
\newcommand{\runicstandardofdismay}{Bannière Runique de Confusion}
\newcommand{\runicstandardofswiftness}{Bannière Runique de Hardiesse}

% Anvil Runes

\newcommand{\runeofstorms}{Rune des Tempêtes}
\newcommand{\runeofshattering}{Rune Fracassante}
\newcommand{\runeofcleansing}{Rune de Purge}

% Characters rules

\newcommand{\monsterseeker}{Tueur de Monstres}
\newcommand{\grimresolve}{Opiniâtre}
\newcommand{\runecraftmastery}{Science Runique}
\newcommand{\mastermarksman}{Maître en Balistique}
\newcommand{\clanleader}{Chef de Clan}
\newcommand{\runecraft}{Artisanat Runique}

% Other units rules

\newcommand{\seenitall}{Radoteurs}
\newcommand{\wallofiron}{Mur d'Acier}
\newcommand{\guildcraftedhandgun}{Arquebuse Manufacturée}
\newcommand{\forgegun}{Fusil de la Forge}
\newcommand{\shrapnelgrenades}{Grenades Shrapnel}
\newcommand{\shrapnelbombs}{Bombes Shrapnel}
\newcommand{\flameculverin}{Couleuvrine Incendiaire}
\newcommand{\whirlingchainsofdoom}{Chaînes du Jugement}
\newcommand{\swiwelgun}{Canon Pierrier}
\newcommand{\majestyofhighkings}{Majesté des Hauts Rois}

% War machines rules

\newcommand{\flamingshot}{Tir Enflammé}
\newcommand{\engineeringrune}{Rune d'Ingénierie}
\newcommand{\organgun}{Canon Orgue}
\newcommand{\organguns}{Canons Orgues}
\newcommand{\runecrafted}{Fonte Runique}
\newcommand{\runestruck}{Frappé de Runes}
\newcommand{\dwarfballista}{Baliste Naine}
\newcommand{\flamecannon}{Canon Incendiaire}
\newcommand{\runicanvil}{Enclume Runique}

%%% Names

% Characters

\newcommand{\king}{Roi Nain}
\newcommand{\daemonseeker}{Chasseur de Démon}
\newcommand{\runicmaster}{Maître-Artisan Runique}
\newcommand{\runicmasters}{Maîtres-Artisans Runique}
\newcommand{\masterengineer}{Maître Ingénieur}
\newcommand{\thane}{Thane}
\newcommand{\dwarfengineer}{Ingénieur Nain}
\newcommand{\runicsmith}{Forgeron Runique}
\newcommand{\runicsmiths}{Forgerons Runiques}
\newcommand{\dragonseeker}{Chasseur de Dragon}

% Core

\newcommand{\clanwarrior}{Guerrier des Clans}
\newcommand{\clanwarriors}{Guerriers des Clans}
\newcommand{\greybeard}{Barbe-Grise}
\newcommand{\greybeards}{Barbes-Grises}
\newcommand{\clanmarksman}{Sentinelle des Clans}
\newcommand{\clanmarksmen}{Sentinelles des Clans}

% Special

\newcommand{\deepwatchSINGULAR}{Garde des Profondeurs}
\newcommand{\deepwatch}{Garde des Profondeurs}
\newcommand{\kingsguard}{Phalange Royale}
\newcommand{\forgewarden}{Veilleur des Forges}
\newcommand{\forgewardens}{Veilleurs des Forges}
\newcommand{\seeker}{Chasseur Funeste}
\newcommand{\seekers}{Chasseurs Funestes}
\newcommand{\steamcopter}{Autogyre à Vapeur}
\newcommand{\steamcopters}{Autogyres à Vapeur}
\newcommand{\miner}{Mineur Nain}
\newcommand{\miners}{Mineurs Nains}
\newcommand{\fieldartillery}{Artillerie de Campagne}

% Rare

\newcommand{\vengeanceseeker}{Chasseur Venge-Serments}
\newcommand{\vengeanceseekerQRS}{Chasseur Venge-Ser.}
\newcommand{\ranger}{Patrouilleur}
\newcommand{\rangers}{Patrouilleurs}
\newcommand{\holdguardian}{Gardien des Forts}
\newcommand{\holdguardians}{Gardiens des Forts}
\newcommand{\grudgebuster}{Brise-Rancunes}
\newcommand{\anvilofpower}{Enclume de Pouvoir}
\newcommand{\anvilsofpower}{Enclumes de Pouvoir}
\newcommand{\grudgemachine}{Artillerie des Rancunes}

% Mounts

\newcommand{\shieldbearers}{Porteurs de Bouclier}
\newcommand{\warthrone}{Trône de Guerre}
\renewcommand{\labels@charactermounts}{Montures du Roi Nain}
\renewcommand{\mountstitle}{%
\def\logolocalpath{../Layout/pics/logo_mount.png}%
\vspace*{1.5cm}\bigtitle{\labels@charactermounts}%
\renewcommand{\profilecategory}{\labels@mounts}%
}

% Short names and multiprofiles names

\newcommand{\pilot}{Pilote}
\newcommand{\copter}{Autogyre}
\newcommand{\attackcopter}{Autogyre d'Assaut}
\newcommand{\steambomber}{Bombardier à Vapeur}
\newcommand{\flamecopter}{Autogyre Incendiaire}
\newcommand{\machine}{Machine}
\newcommand{\crew}{Servant}
\newcommand{\grudgebustercrew}{Équipage}
\newcommand{\anvil}{Enclume}
\newcommand{\anvilcrew}{Garde de l'Enclume}

% Profile wording

% Runes

\newcommand{\maytakearunicweapon}{Arme Runique}
\newcommand{\maytakerunicitems}{Objets Runiques}
\newcommand{\battleruneschoice}[1]{Jusqu'à #1 Runes de Batailles différentes}
\renewcommand{\perrune}{/rune}

% Others

\newcommand{\otherworldlymodels}{Figurines \otherworldly{}}
\newcommand{\onlyonemaytakevanguard}{Une seule unité de \clanwarriors{} ou de \greybeards{} peut gagner la règle \vanguard{}}
\newcommand{\mayswapcrossbowfor}{\optionschoice{Échanger l'\crossbow{} contre :}}
\newcommand{\seekersskirmishersoption}{\skirmishers{} (\oneofakind{}, 15 fig. max.)}
\newcommand{\copterchoice}{Doit choisir l'une des trois unités :}
\newcommand{\mayaddonemodelfor}{Peut ajouter une figurine pour}
\newcommand{\copterskirmishoption}{\skirmishers{} (si 2 fig.)}

% Profile Rules

% Characters

\newcommand{\grimresolverule}{%
La figurine gagne une attaque supplémentaire par figurine ennemie en contact socle à socle avec elle.
}

\newcommand{\monsterseekerrule}{%
La figurine gagne la règle \multiplewounds{2}{\monstrousbeasts{}, \monstrouscavalry{}, \monstrousinfantry{}, \monster{} et \riddenmonster{}}.
}

\newcommand{\runecraftmasteryrule}{%
L'\runicmaster{} et son unité disposent de la règle \armourpiercing{1}. Il peut dissiper comme un Maître Sorcier et utiliser ses Runes de Bataille comme des Objets de Sort, avec au choix :
\begin{itemize}[label={-}]
\item Niveau de Puissance 4, \range{6}.
\item Niveau de Puissance 5, \range{12}.
\end{itemize}
}

\newcommand{\mastermarksmanrule}{%
Au lieu d'utiliser la règle \engineer{}, un \masterengineer{} peut, au début de la Phase de Tir et pour la durée de celle-ci, octroyer un des bonus suivants à une unité d'Infanterie alliée située à moins de \distance{6} :
\begin{itemize}[label={-}]
\item L'unité gagne un bonus de +1 pour toucher au Tir.
\item L'unité peut relancer les jets pour blesser ayant donné \result{1} au Tir.
\item La portée des Armes de Tir de l'unité augmente de \distance{1D6+1}.
\end{itemize}

Un bonus donné ne peut pas se cumuler sur une même unité.
}

\newcommand{\clanleaderrule}{%
Le \thane{} et tous les \clanwarriors{} de son unité gagnent la règle \fightinextrarank{}.
}

\newcommand{\runecraftrule}{%
Le \runicsmith{} et son unité disposent de la règle \armourpiercing{1}. Il peut dissiper comme un Sorcier Apprenti et utiliser ses Runes de Bataille comme des Objets de Sort, avec au choix :
\begin{itemize}[label={-}]
\item Niveau de Puissance 4, Unité du Lanceur.
\item Niveau de Puissance 5, \range{6}.
\end{itemize}
}

% Anvil

\newcommand{\runicanvilrule}{%
L'\anvilofpower{} peut dissiper comme un Sorcier Apprenti. Elle peut choisir jusqu'à trois Runes, gratuitement parmi les Runes de Bataille, et à leur prix parmi les Runes de l'Enclume ci-dessous. Toutes les Runes sont des Objets de Sort avec une portée de \distance{36}. L'Enclume peut lancer les Runes de Bataille avec un Niveau de Puissance de 4 et les Runes de l'Enclume avec un Niveau de Puissance de 5.
}

\newcommand{\anvilrunes}{%
\renewcommand{\arraystretch}{3.2}
\begin{tabular}{@{}>{\raggedright}m{3cm}M{2.2cm}M{2.2cm}m{7.8cm}@{}}
\textbf{Nom} &
\textbf{Type} &
\textbf{Durée} &
\centering\textbf{Effet}
\tabularnewline\hline

\runeofstorms\newline (\pts{5}) &
Dégâts\newline Malédiction &
Immédiat\newline Dure un Tour &
La cible subit immédiatement 1D6 touches de Force 6 avec la règle \lightningattacks{}. La cible ne peut plus utiliser la règle \fly{} pendant la durée du sort.
\tabularnewline\hline

\runeofshattering\newline (\pts{5}) &
Dégâts\newline Malédiction &
Immédiat\newline Dure un Tour &
La cible subit immédiatement 2D6 touches de Force 4. Toute unité subissant au moins une blessure de ces touches subit un malus de -1 sur ses jets pour toucher au Corps à Corps, et traite tous les décors, y compris les Terrains Dégagés, comme des Terrains Dangereux (1).
\tabularnewline\hline

\runeofcleansing\newline (\free{}) &
Universelle &
Immédiat &
Choisissez un sort de type Dure un tour ou Reste en jeu affectant la cible. Ce sort prend immédiatement fin. 
\tabularnewline
\end{tabular}
\renewcommand{\arraystretch}{1.2} % Return to default
}

% Warmachines

\newcommand{\flamingshotrule}{%
Les tirs gagnent les règles \flamingattacks{} et \magicalattacks{}.
}

\newcommand{\engineeringrunerule}{%
Ajoutez +4 aux jets de dés sur la table des Incidents de Tir.
}

\newcommand{\catapultrule}{%
\textbf{\artilleryweapon} de type \textbf{\catapult{} (\distance{3})}.\newline
\range{12-60}, \Strength{} 3 [9], [\multiplewounds{\ordnance}{}].
}

\newcommand{\catapultrunecraftedrule}{%
Toutes\newline les touches de la \catapult{} gagnent +1 en Force et les règles \magicalattacks{} et \armourpiercing{1}.
}

\newcommand{\cannonrule}{%
Peut tirer de deux manières :
\begin{customsubitemize}
\item Comme une \textbf{\artilleryweapon} de type \textbf{\cannon{} (\distance{1D6})}.\newline
\range{60}, \Strength{} 10, \multiplewounds{\ordnance}{}.
\item Comme une \textbf{\artilleryweapon} de type \textbf{\volleygun}.\newline
\range{12}, \Strength{} 4, \armourpiercing{3}, \multipleshots{2D6}.
\end{customsubitemize}
}

\newcommand{\cannonrunecraftedrule}{%
Toutes\newline les touches du \cannon{} gagnent +1 en Force et la règle \magicalattacks{}.
}

\newcommand{\organgunrule}{%
\textbf{\artilleryweapon} de type \textbf{\volleygun}.\newline
\range{30}, \Strength{} 5, \armourpiercing{1}, \multipleshots{2D6*2}.
}

\newcommand{\organgunrunecraftedrule}{%
Toutes\newline les touches du \organgun{} gagnent un bonus de +1 sur les jets pour blesser et la règle \magicalattacks{}.
}

\newcommand{\flamecannonrule}{%
\textbf{\artilleryweapon} de type \textbf{\flamethrower}.\newline
\range{12}, \Strength{} 5, \flamingattacks{}, \multiplewounds{1D3}{}.
}

\newcommand{\flamecannonrunecraftedrule}{%
Toutes les touches du \flamecannon{} gagnent la règle \magicalattacks{}.
}

\newcommand{\runestruckrule}{%
Le \flamecannon{} gagne \distance{+3} en portée et la règle \magicalattacks{}.
}

\newcommand{\dwarfballistarule}{%
\textbf{\artilleryweapon} de type \textbf{\boltthrower}.\newline
\range{48}, \Strength{} 6, \multiplewounds{1D3}{}, \armourpiercing{6}.
}

\newcommand{\dwarfballistarunecraftedrule}{%
Toutes les touches de la \dwarfballista{} gagnent un bonus de +1 pour toucher contre les unités volantes et la règle \magicalattacks{}.
}

\newcommand{\grudgemachinenote}{%
L'\grudgemachine{} est \oneofakind{} si l'armée contient au moins 3 figurines d'\fieldartillery{}. Ce chiffre est réduit à 2 pour les Patrouilles et augmenté à 6 pour les Grandes Armées. Ignorez cette restriction si toutes les figurines d'\grudgemachine{} prennent l'amélioration \dwarfballista{}.
}

% Other units

\newcommand{\seenitallrule}{%
Toute unité alliée à moins de \distance{6} peut relancer ses tests de Panique ratés.
}

\newcommand{\wallofironrule}{%
La \wardsave{} accordée par la règle \shieldwall{} est toujours de 5+.
}

\newcommand{\majestyofhighkingsrule}{%
Si un Général est monté sur le \warthrone{}, la portée de sa règle \inspiringpresence{} passe à \distance{18}. Toutes les unités amies à moins de \distance{6} gagnent la règle \swiftstride{}. Si le Personnage sur le \warthrone{} était dans une unité lorsqu'il est tué, l'unité gagne la règle \frenzy{}.

\vspace*{5pt}L'\fieldartillery{} passe des Unités Spéciales aux Unités Rares.
}

% Weapons

\newcommand{\guildcraftedhandgunrule}{%
Arme de Tir. Type : Arquebuse. Accorde un bonus de +1 pour toucher au tir. Une unité équipée de cette arme ne peut pas porter d'\gw{} ni avoir de Porte-étendard Vétéran.
}

\newcommand{\forgegunrule}{%
Arme de Tir. \range{18}, Force 5, \flamingattacks{}, \quicktofire{}.
}

\newcommand{\shrapnelgrenadesrule}{%
\sweepingattack{} qui peut être utilisée une fois par partie. L'unité ennemie subit 1D3 touches de Force 3 pour chaque \copter{} dans l'unité.
}

\newcommand{\shrapnelbombsrule}{%
\sweepingattack{}. L'unité ennemie subit 1D6*2 touches de Force 3 avec la règle \armourpiercing{1}. Si un \result{6} naturel est obtenu sur le dé du nombre de touches, les \shrapnelbombs{} ne pourront plus être utilisées de la partie après cette attaque.
}

\newcommand{\flameculverinrule}{%
\textbf{\artilleryweapon} de type \textbf{\flamethrower}.\newline
\range{3}, Force 3, \flamingattacks{}. Ne peut pas tirer si la figurine a effectué une Marche Forcée à ce Tour de Joueur. Ignorez le malus de -1 sur la table d'Incident de Tir.
}

\newcommand{\whirlingchainsofdoomrule}{%
Arme de Corps à Corps. Le porteur gagne +1 en Force, la règle \randomattacks{3D3} et attaque à Initiative 10. Il ne peut pas être rejoint par un Personnage.
}

\newcommand{\swiwelgunrule}{%
Arme de Tir. \range{24}, Force 4, \armourpiercing{1}, \multipleshots{4}, \quicktofire{}.
}







\begin{document}

\newgeometry{margin=1in}

% Table options
\arrayrulecolor{black!30}
\setlength{\arrayrulewidth}{2pt}
\renewcommand{\arraystretch}{1.2}

\begin{titlepage}
\begin{center}

\ifdef{\booktitle}{}{\newcommand{\booktitle}{Missing title}}
\ifdef{\version}{}{\newcommand{\version}{Missing version}}

{\antiquefont\fontsize{40}{48}\selectfont\noindent\labels@fantasybattles

\labels@NinthAge}

\vspace*{0.5cm}
\ifdef{\booklogo}{%
\includegraphics[height=10cm]{\booklogo}%
}{%
\includegraphics[height=10cm]{../Layout/pics/logo_9th.png}%
}

\vspace*{-1cm}
{\antiquefont\fontsize{50}{60}\selectfont \booktitle
\vspace{0.4cm}

\fontsize{14}{16.8}\selectfont \labels@armyrules{}

Beta v\version{} - \today{}}

\ifdef{\frenchversion}{{\fontsize{14}{16.8}\selectfont \vspace{0.2cm}\noindent\texttt{VF \frenchversion}}}{}
\vfill

\begin{tabular}{@{}m{2cm}@{\hskip 20pt}m{13cm}@{}}
\includegraphics[width=2cm]{../Layout/pics/seal_9th.png} &
{\fontsize{10}{12}\selectfont \textcolor{black!50}{\noindent\labels@frontpagecredits}}

\ifdef{\frontpageaddstuff}{{\fontsize{10}{12}\selectfont \noindent\textcolor{black!50}{\frontpageaddstuff}}}{}

\vspace*{10pt}
\noindent{\fontsize{10}{12}\selectfont \textcolor{black!50}{\labels@license}}
\tabularnewline
\end{tabular}


\end{center}

\newpage

\thispagestyle{empty}

{\fontsize{10}{12}\selectfont

\ifdef{\labels@introduction}{\labels@introduction}{\vphantom{1pt}}
\vfill

\noindent\newrule{\labels@rulechanges}

\bigskip
\noindent \labels@latexcredit
}


\end{titlepage}

\restoregeometry

\startarmywiderules

\armywideruleentry{\hewnoutofmountains}

Tant qu'il reste sur la table au moins une unité alliée d'une armée des Forteresses Naines, tous les sorts lancés par les figurines ennemies voient leur valeur de lancement, ou leur niveau de puissance pour les Objets de Sorts, augmentée de +1. La valeur de lancement des sorts restant en jeu ne change pas pour ce qui est de leur dissipation.

\armywideruleentry{\ancientgrudge}

Une armée des Forteresses Naines a un nombre de \ancientgrudges{} défini comme suit :
\begin{itemize}[label={-}]
\item Une \ancientgrudge{} par \king{} présent dans l'armée.
\item 1D3 \ancientgrudges{} additionnelles pour chaque \king{} monté sur un \warthrone{}.
\item Une \ancientgrudge{} additionnelle si l'armée combat des figurines issues des livres d'armées Peaux Vertes ou Marée de Vermine.
\end{itemize}

Après les déploiements des deux joueurs, vous pouvez désigner une figurine ou une unité ennemie pour chaque \ancientgrudge{}. Aucune figurine ou unité ne peut être ciblée plus d'une fois. Toutes les unités de votre armée gagnent la règle \hatred{} contre ces figurines ou unités.

Si c'est une unité qui est choisie, toute figurine qui la rejoint subit aussi les effets de la \ancientgrudge{} tant qu'elle reste dans l'unité. Un Personnage qui quitte une unité contre laquelle il existe une \ancientgrudge{} n'en subit donc plus les effets. Si une figurine individuelle est désignée et qu'elle rejoint une unité, les effets de la \ancientgrudge{} ne s'étendent pas à l'unité entière. 

\closearmywiderules



\vspace*{1.5cm}
\startarmyspecialrules

\armyspecialruleentry{\relentless}

Les unités d'Infanterie contenant uniquement des figurines dotées de cette règle peuvent tripler leur mouvement lors d'une Marche Forcée au lieu de le doubler. Ceci est étendu aux situations où le mouvement est d'ordinaire limité au double de la valeur de mouvement de la figurine, telles que les Roues, les Reformations, les mouvements d'un Personnage au sein d'une unité, etc. La valeur de mouvement est alors également triplée.

\armyspecialruleentry{\sturdy}

La figurine gagne la règle \thunderouscharge{}. De plus, elle ne souffre pas de la pénalité de -1 pour toucher pour Tenir sa Position et Tirer en réaction à une charge. 

\armyspecialruleentry{\shieldwall}

Lorsque la figurine utilise un Bouclier, elle gagne une \wardsave{6} contre les attaques de Corps à Corps non Spéciales provenant d'unités ennemies engagées sur le front. Cette \wardsave{} passe à 5+ durant chaque première Manche de Corps à Corps pendant laquelle son unité a été chargée. 

\armyspecialruleentry{\entrench}

Une fois le déploiement terminé mais avant de placer les Éclaireurs, la figurine peut retrancher une seule Machine de Guerre qui comptera comme étant à Couvert Lourd. Si la Machine de Guerre bouge pour une raison ou une autre, elle perd définitivement ce couvert.

\armyspecialruleentry{\holdstone}

Au début de n'importe quelle Manche de Corps à Corps, la figurine peut déclarer qu'elle l'utilise sa \holdstone{}. A partir de ce moment, l'unité du porteur considère ses flancs et son arrière comme étant de front en ce qui concerne les Attaques de Soutien et les règles \parry{} et \shieldwall{}. Elle peut néanmoins toujours être Désorganisée. Tant que la \holdstone{} est active, le porteur doit toujours lancer et relever des défis, à moins qu'une autre figurine amie ne le fasse avant. L'unité ne peut pas poursuivre directement après le combat où la \holdstone{} a été utilisée.

Les effets de la \holdstone{} disparaissent immédiatement une fois le combat terminé ou à la fin de la phase durant laquelle son porteur est tué.

Tant que le porteur est en vie, la Pierre peut être activée à nouveau lors des combats suivants.

\armyspecialruleentry{\thebiggertheyare}

Les Attaques de Corps à Corps de la figurine blessent toujours sur 4+ ou mieux. Elle gagne également la règle \swiftstride{} lorsqu'elle fait le jet d'une charge vers une unité contenant au moins une figurine d'un des types suivant : Bête Monstrueuse, Cavalerie Monstrueuse, Infanterie Monstrueuse, Monstre ou Monstre Monté.

\armyspecialruleentry{\comewithme}

Juste avant qu'elle soit retirée comme perte suite à une attaque de Corps à Corps durant la Phase de Corps à Corps, la figurine peut porter immédiatement une ultime attaque de Corps à Corps avec l'arme utilisée pendant cette manche. L'attaque doit être portée au choix sur la figurine adverse l'ayant achevé ou bien son unité, auquel cas cette attaque est distribuée comme un tir. Les bonus et règles spéciales sont appliquées normalement, mais la figurine ne peut porter qu'une seule attaque au plus dont la Force ne peut pas dépasser 5. Cette règle ne peut pas être utilisée suite à des pertes infligées par des \impacthits{}. Une unité ne peut utiliser cette règle que si elle a au moins autant de rangs que de colonnes au début de la Manche de Corps à Corps.

\closearmyspecialrules






\vspace*{1.5cm}
\startarmyarmoury

\startitemlistonecol

\listitemonecol{\wyrmslayerrocket} Arme de Tir. \range{24}, Force 6, \flamingattacks{}, \multiplewounds{1D3}{}, \reload{}.

\listitemonecol{\forgerepeater} Arme de Tir. \range{18}, Force 5, \flamingattacks{}, \multipleshots{4}, \quicktofire{}.

\enditemlistonecol

\closearmyarmoury





\toctarget{battlerunestitle}{\startarmynewsection{Runes de Bataille}}\subtitle{}

Les \runicmasters{}, les \runicsmiths{} et les \anvilsofpower{} peuvent disposer de ces Runes, en suivant les règles détaillées dans leur profil, et les lancer comme des Objets de Sort. Aucune Rune de Bataille ne peut être choisie plus de deux fois dans une même armée. Ce chiffre est réduit à un pour les Patrouilles et augmenté à quatre pour les Grandes Armées.

\vspace*{1.5cm}
\def\Z{\vphantom{\parbox[c]{2.5cm}{\fontsize{25}{30}\selectfont Some\newline Very\newline Long}}}

\newcommand{\runewidth}{0.8cm}
\begin{tabular}{@{}c@{}c@{\hskip 8pt}>{\raggedright}m{3.3cm}M{2.1cm}M{2.1cm}m{7.4cm}@{}}
&&
\textbf{Nom} &
\textbf{Type} &
\textbf{Durée} &
\centering\textbf{Effet}
\tabularnewline\hline

\Z{} & \includegraphics[valign=m, width=\runewidth]{pics/battlerune_metal.png} &
\runeofmetal &
Amélioration &
Dure un Tour &
La cible peut relancer ses jets de Sauvegarde d'Armure ratés.
\tabularnewline\hline

\Z{} & \includegraphics[valign=m, width=\runewidth]{pics/battlerune_resolve.png} &
\runeofresolve &
Amélioration &
Immédiat &
L'unité alliée ciblée peut effectuer un Mouvement Magique de \distance{6}. 
\tabularnewline\hline

\Z{} & \includegraphics[valign=m, width=\runewidth]{pics/battlerune_reckoning.png} &
\runeofreckoning &
Amélioration &
Dure un Tour &
L'unité alliée ciblée peut relancer ses jets pour toucher ratés au Corps à Corps.
\tabularnewline\hline

\Z{} & \includegraphics[valign=m, width=\runewidth]{pics/battlerune_oaths.png} &
\runeofoaths &
Amélioration &
Dure un Tour &
La cible gagne les règles \stubborn{} et \immunetopsychology{}.
\tabularnewline\hline

\Z{} & \includegraphics[valign=m, width=\runewidth]{pics/battlerune_gleaming.png} &
\runeofgleaming &
Amélioration &
Dure un Tour &
La cible gagne les règles \hardtarget{} et \distracting{}.
\tabularnewline\hline

\Z{} & \includegraphics[valign=m, width=\runewidth]{pics/battlerune_resilience.png} &
\runeofresilience &
Amélioration &
Dure un Tour &
Les jets pour blesser contre la cible subissent un malus de -1. Plusieurs instances de ce sort sur une même unité ne s'additionnent pas.
\tabularnewline

\end{tabular}

\closearmynewsection




\toctarget{dwarvenrunestitle}{\startarmynewsection{Runes Naines}}\subtitle{}

Les armées des Forteresses Naines n'utilisent pas les Objets Magiques Communs, sauf les Bannières Magiques, mais créent plutôt les leurs en utilisant des Runes Naines. Les objets ainsi forgés sont appelés Objets Runiques et suivent les règles habituelles des Objets Magiques. Une Arme Runique compte donc par exemple comme une Arme Magique. Chaque combinaison de Runes correspond à un unique Objet Magique qui ne peut pas être dupliqué au sein d'une même armée. Une Rune peut être dupliquée sur un même objet s'il n'y a pas de mention contraire.

\begin{multicols}{2}\raggedcolumns

\subtitle{Runes d'Arme}\vspace{5pt}

Toute \hw{} ou \pw{} peut être enchâssée d'un maximum de trois Runes parmi cet ensemble.

\startpricelist

\pricelistitem{\runeofpenetrating}{40/30}\oneofakind{}.

Les attaques portées avec l'arme gagnent la règle\armourpiercing{6}.

\pricelistitem{\runeofdestruction}{40/20}\oneofakind{}.

Les attaques portées avec l'arme gagnent la règle \multiplewounds{1D3}{}.

\pricelistitem{\runeofsmashing}{35/25}\oneofakind{}.

Les attaques portées avec l'arme ont Force 10 quand la cible possède une Endurance de 5 ou plus.

\pricelistitem{\runeofmight}{20}Les attaques portées avec l'arme ont un bonus de +1 en Force pour chaque Rune de Puissance enchâssée.

\pricelistitem{\runeofprecision}{20}Le porteur gagne la règle \lightningreflexes{}.

\pricelistitem{\runeofquickening}{20}Le porteur attaque avec un bonus de +2 en Initiative pour chaque Rune d'Accélération enchâssée.

\pricelistitem{\runeofcraftsmanship}{20/15}\oneofakind{}.

L'arme devient une \gw{}.

\pricelistitem{\runeoffury}{15}Le porteur gagne +1 Attaque par Rune de Fureur enchâssée.

\pricelistitem{\runeoffire}{5}Les attaques portées avec l'arme ont la règle \flamingattacks{}.

\endpricelist

\columnbreak
\subtitle{Runes d'Armure}\vspace{5pt}

Tout ensemble d'armure peut être enchâssé d'un maximum de trois Runes parmi cet ensemble.

\startpricelist

\pricelistitem{\runeofresistance}{45}\oneofakind{}.

Les jets pour blesser réussis contre le porteur doivent être relancés.

\pricelistitem{\runeofbronze}{30}Le porteur gagne +1 Point de Vie.

\pricelistitem{\runeofsteel}{30}Le porteur peut relancer ses jets de Sauvegarde d'Armure ratés.

\pricelistitem{\runeofmountains}{30}Le porteur gagne +1 en Endurance jusqu'à un maximum de 6.

\pricelistitem{\runeofaegis}{15}Le porteur gagne la règle \magicresistance{1} ou augmente cet effet de +1 par Rune, jusqu'à \magicresistance{4} au mieux. Cela n'augmente pas la \magicresistance{} qui serait obtenue par le biais d'une autre figurine de l'unité du porteur.

\pricelistitem{\runeofiron}{10} Le porteur gagne une \innatedefence{6}. Une autre Rune peut être ajoutée pour passer cet effet à (5+).

\pricelistitem{\runeofcrushing}{10} Le porteur gagne la règle \impacthits{+1}.

\endpricelist

\columnbreak
\subtitle{Runes de Talisman}\vspace{5pt}

Un maximum de 3 Runes de cet ensemble peuvent être sélectionnées et comptent alors comme un unique Talisman.

\startpricelist

\pricelistitem{\runeofretribution}{35}\oneperarmy{}.

Usage Unique. Peut être activée au début de n'importe quelle Phase de Mouvement. Les unités amies à moins de \distance{6} du porteur gagnent la règle \devastatingcharge{} jusqu'à la fin du Tour de Joueur.

\pricelistitem{\runeofdragonsbreath}{35}Le porteur gagne la règle \breathweapon{Force 4, \flamingattacks{} et \magicalattacks{}}.

\pricelistitem{\runeofcourage}{30}Usage Unique. Peut être activée au début de n'importe quelle Manche de Corps à Corps. Le porteur gagne la règle \stubborn{} pour la durée de la Phase.

\pricelistitem{\runeofinfamy}{20}Le porteur provoque la \fear{}.

\pricelistitem{\runeofshielding}{15}Le porteur gagne une \wardsave{6} ou améliore la valeur de cette \wardsave{} de 1 par Rune. L'effet de cette Rune ne s'additionne \textbf{qu}'avec lui-même et la règle \shieldwall{}, et jusqu'à obtenir 4+ au mieux. Ceci n'affecte pas l'interaction avec la \magicresistance{}.

\pricelistitem{\runeofkinship}{10}Le porteur gagne les règles \vanguard{} et \scout{} mais ne peut pas être monté sur un \warthrone{}.

\pricelistitem{\runeoftheforge}{5}Le porteur gagne la règle \fireborn{}.

\endpricelist

\subtitle{Runes Cabalistiques}\vspace{5pt}

Les figurines avec la règle \runecraft{} ou \runecraftmastery{} peuvent prendre un maximum de 3 Runes de cet ensemble qui comptent alors comme un unique Objet Cabalistique.

\startpricelistNSP

\pricelistitem{\runeofgrounding}{50}\oneofakind{}.

À la fin de votre Phase de Magie, vous pouvez mettre de côté un Dé de Magie et l'ajouter à votre réserve à la Phase de Magie suivante, immédiatement après avoir jeté les dés pour les Flux de Magie.

\pricelistitem{\runeofmastery}{40}Une fois par Phase de Magie, le porteur peut relancer un Dé de Pouvoir utilisé pour un Objet de Sort.

\columnbreak
\vspace*{1.93cm}

\pricelistitem{\runeofdevouring}{35}\oneperarmy{}.

Usage Unique. Le joueur peut décider d'utiliser cette Rune plutôt que de tenter une dissipation. Le sort est lancé normalement, mais il est ensuite perdu et ne peut plus être relancé de la partie. Cette Rune ne peut pas être combinée avec une Rune de Déni sur un même objet.

\pricelistitem{\runeofdenial}{35}\oneperarmy{}.

Usage Unique. Le joueur peut décider d'utiliser cette Rune plutôt que de tenter une dissipation. Le sort est automatiquement dissipé.

\pricelistitem{\runeofbrilliance}{35} Le porteur ajoute un bonus de +1 à tous ses jets de dissipation ainsi qu'aux jets de lancement de sorts liés à des Objets de Sort. Cela passe outre les restrictions habituelles empêchant les modificateurs de lancement pour les Objets de Sort.

\pricelistitem{\runeofharnessing}{25} Le joueur profite d'un bonus de +1 sur ses jets de \channel{} tandis que l'adversaire subit un malus de -1 sur les siens.

\endpricelistNSP

\subtitle{Bannières Runiques}\vspace{5pt}

Les armées des Forteresses Naines peuvent aussi bien utiliser les Bannières Magiques Communes que les Bannières Runiques listées ci-dessous. Un Porteur de la Grande Bannière peut prendre une Bannière Magique Commune ou une Bannière Runique dans son budget d'Objets Runiques.

\startpricelist

\pricelistitem{\runicstandardofshielding}{45}Toutes les unités alliées à moins de \distance{6} du porteur gagnent une \wardsave{5} contre les Attaques de Tir.

\pricelistitem{\runicstandardofdismay}{30}Les unités chargeant l'unité du porteur réduisent leur distance de charge de \distance{1D6}. Jetez le dé pour chaque unité qui charge avant de les déplacer. Cela peut provoquer une charge ratée.

\pricelistitem{\runicstandardofswiftness}{25}L'unité du porteur gagne la règle \vanguard{}.

\endpricelist

\end{multicols}
\closearmynewsection








%%% START OF THE ARMYLIST - Translators shouldn't have to edit it %%%


%%% v0.99.9

\armylist

\lordstitle

\showunit{
	name={\king},
	cost=135,
	profile={< 3 7 4 4 5 3 4 4 10},
	type=\infantry{},
	unitsize=1,
	basesize=20x20,
	commontype=\dwarvencommonrules{},
	commonspecialrules={\relentless{},\shieldwall{},\sturdy},
	armour={\platearmour{}},
	options={
		\maytakerunicitems{}=\upto{}<125,
		\holdstone{}=30,
		\shield{}=10,
		\shootingweapononechoice{
			\pistol{}=4,
			\crossbow{}=8,
			\handgun{}=8,
		},
		\gw{}=10,
	},
	mounts={\shieldbearers{}=50,\warthrone{}=150},
}

\mountstitle

\showunit{
	name={\shieldbearers},
	profile={ < 3 5 3 4 4 4 2 2 10},
	type=\infantry{},
	basesize=40x20,
	commontype=\dwarvencommonrules{},
	commonspecialrules={\relentless{},\sturdy{}},
	armour={\mountsprotection{5}},
}

\showunit{
	name={\warthrone},
	profile={ < 3 5 3 4 4 6 2 4 10},	
	type=\infantry{},
	basesize=40x60,
	commontype=\dwarvencommonrules{},
	commonspecialrules={\relentless{},\sturdy{}},
	specialrules={\hatred{},\majestyofhighkings{}}, 
	armour={\mountsprotection{5}},
	unitrules={
		\unitrule{\majestyofhighkings{}}{\majestyofhighkingsrule{}}
	},
}	




\clearpage
\def\logolocalpath{../Layout/pics/logo_lord.png}%
\renewcommand{\profilecategory}{\labels@lords}%

\showunit{
	name={\daemonseeker},
	cost={130},
	profile={< 3 7 4 5 5 3 5 5 10},
	type=\infantry{},
	unitsize=1,
	basesize=20x20,
	commontype=\dwarvencommonrules{},
	commonspecialrules={\relentless{},\sturdy{}},
	commontypeB=\seekerscommonrules{},
	commonspecialrulesB={\unbreakable{},\weaponmaster{},\wardsave{6},\thebiggertheyare{},\comewithme{}},
	specialrules={\vanguard{},\lethalstrike{},\hatred{} (\otherworldlymodels{}),\notaleader{}},
	weapons={\pw{},\gw{}},
	options={
		\maytakearunicweapon{}=\upto{}<125,
		\uptotwoofthefollowing{
			\magicresistance{1}=15,
			\fear{}=20,
			\monsterseeker{}=30,
			\grimresolve{}=40,
		},
	},
	unitrules={
		\unitrule{\monsterseeker{}}{\monsterseekerrule{}}
		\unitrule{\grimresolve{}}{\grimresolverule{}}
	},
}	

\showunit{
	name={\runicmaster},
	cost={125},
	profile={< 3 6 4 4 5 3 3 2 9},
	type= \infantry{} ,
	basesize=20x20,
	unitsize=1,
	commontype=\dwarvencommonrules{},
	commonspecialrules={\relentless{},\shieldwall{},\sturdy{}},
	specialrules={\channel{},\magicresistance{2},\runecraftmastery{}},
	armour={\platearmour{}},
	options={
		\maytakerunicitems{}=\upto{}<125,			
		\shield{}=3,
		\gw{}=8,
		\battleruneschoice{4}=\perrune{}<5,
	},
	unitrules={
		\unitrule{\runecraftmastery{}}{\runecraftmasteryrule{}}
	},
}

\showunit{
	name={\masterengineer},
	cost={100},
	profile={< 3 6 4 4 5 3 3 2 9},
	type=\infantry{},
	unitsize=1,
	basesize=20x20,
	commontype=\dwarvencommonrules{},
	commonspecialrules={\relentless{},\shieldwall{},\sturdy{}},
	specialrules={\engineer{},\mastermarksman{},\entrench{}},
	armour={\platearmour{}},
	options={
		\maytakerunicitems{}=\upto{}<100,
		\shield{}=3,
		\shootingweapononechoice{
			\pistol{}=4,
			\braceofpistols{}=5, 
			\crossbow{}=8,
			\handgun{}=8,
			\forgerepeater{}=10,
			\wyrmslayerrocket{}=20,
		},
		\gw{}=8,
	},
	additional={
		\def\tempunitrules{\unitrule{\mastermarksman{}}{\mastermarksmanrule{}}}
		\unitrules{\tempunitrules}
	},
}	
	






	

\heroestitle

\showunit{
	name={\thane},
	cost={70},
	profile={< 3 6 4 4 5 2 3 3 9},
	type=\infantry{},
	unitsize=1,
	basesize=20x20,
	commontype=\dwarvencommonrules{},
	commonspecialrules={\relentless{},\shieldwall{},\sturdy{}},
	armour={\platearmour{}},
	options={
		\maytakerunicitems{}=\upto{}<75,
		\holdstone{}=25,
		\onechoiceonly{
			\bsb{}=25,
			\clanleader{}=30,
		},
		\shield{}=8,
		\shootingweapononechoice{
			\pistol{}=4,
			\crossbow{}=8,
			\handgun{}=8,
		},
		\gw{}=10,
	},
	unitrules={
		\unitrule{\clanleader{}}{\clanleaderrule{}}
	},
}

\showunit{
	name={\dwarfengineer},
	cost={65},
	profile={< 3 5 4 4 4 2 3 2 9},
	type=\infantry{},
	unitsize=1,
	basesize=20x20,
	commontype=\dwarvencommonrules{},
	commonspecialrules={\relentless{},\shieldwall{},\sturdy{}},
	specialrules={\entrench{},\engineer{}},
	armour={\platearmour{}},
	options={
		\maytakerunicitems{}=\upto{}<50,
		\shield{}=2,
		\shootingweapononechoice{
			\pistol{}=4,
			\braceofpistols{}=5, 
			\crossbow{}=8,
			\handgun{}=8,
			\forgerepeater{}=10,
			\wyrmslayerrocket{}=20,
		},
		\gw{}=6,
	},
}

\showunit{
	name={\runicsmith},
	cost={70},
	profile={< 3 5 3 4 4 2 3 2 9},
	type=\infantry{},
	unitsize=1,
	basesize=20x20,
	commontype=\dwarvencommonrules{},
	commonspecialrules={\relentless{},\shieldwall{},\sturdy{}},
	specialrules={\runecraft{},\channel{},\magicresistance{1}},
	armour={\platearmour{}},
	options={
		\maytakerunicitems{}=\upto{}<75,			
		\shield{}=2,
		\gw{}=6,
		\battleruneschoice{2}=\perrune{}<5,
	},
	unitrules={
		\unitrule{\runecraft{}}{\runecraftrule{}}
	},
}

\showunit{
	name={\dragonseeker},
	cost={50},
	profile={< 3 6 3 4 5 2 4 4 10},
	type=\infantry{},
	unitsize=1,
	basesize=20x20,
	commontype=\dwarvencommonrules{},
	commonspecialrules={\relentless{},\sturdy{}},
	commontypeB=\seekerscommonrules{},
	commonspecialrulesB={\unbreakable{},\weaponmaster{},\wardsave{6},\thebiggertheyare{},\comewithme{}},
	specialrules={\vanguard{},\lethalstrike{},\notaleader{}},
	weapons={\gw{},\pw{}},
	options={
		\maytakearunicweapon{}=\upto{}<75,
		\uptotwoofthefollowing{
			\magicresistance{1}=15,
			\fear{}=20,
			\monsterseeker{}=20,
			\grimresolve{}=25,
		},
	},
	unitrules={
		\unitrule{\monsterseeker{}}{\monsterseekerrule{}}
		\unitrule{\grimresolve{}}{\grimresolverule{}}
	},
}






\coreunitstitle

\showunit{
	name={\clanwarriors},
	QRSname={\clanwarrior},
	cost={60},
	profile={< 3 4 3 3 4 1 2 1 9},
	costpermodel=7,
	maxmodels=40,
	type=\infantry{},
	unitsize=10,
	basesize=20x20,
	commontype=\dwarvencommonrules{},
	commonspecialrules={\relentless{},\shieldwall{},\sturdy{}},
	armour={\ha{}},
	options={
		\onlyonemaytakevanguard{} \Xmodelsorless{30}=\permodel{}<1,
		\shield{}=\permodel{}<1,
		\throwingweapons{}=\permodel{}<1,
		\weapononechoice{
			\pw{}=\permodel{}<1,
			\spear{}=\permodel{}<1,
			\gw{}=\permodel{}<3,
		},
	},
	commandgroup={champion=10, musician=10, banner=10, veteranstandardbearer=yessir},
}

\showunit{
	name={\greybeards},
	QRSname={\greybeard},
	cost={90},
	profile={< 3 5 3 4 4 1 2 1 9},
	costpermodel=11,
	maxmodels=30,
	type=\infantry{},
	unitsize=10,
	basesize=20x20,
	commontype=\dwarvencommonrules{},
	commonspecialrules={\relentless{},\shieldwall{},\sturdy{}},
	specialrules={\seenitall{},\immunetopsychology{}},
	armour={\ha{}},
	options={
		\onlyonemaytakevanguard{} \Xmodelsorless{20}=\permodel{}<1,
		\shield{}=\permodel{}<1,
		\throwingweapons{}=\permodel{}<2,
		\gw{}=\permodel{}<3,
	},
	unitrules={
		\unitrule{\seenitall{}}{\seenitallrule{}}
	},
	commandgroup={champion=10, musician=10, banner=10, veteranstandardbearer=yessir},
}

\showunit{
	name={\clanmarksmen},
	QRSname={\clanmarksman},
	cost={110},
	profile={ < 3 4 3 3 4 1 2 1 9},
	costpermodel=11,
	maxmodels=25,
	type=\infantry{},
	unitsize=10,
	basesize=20x20,
	commontype=\dwarvencommonrules{},
	commonspecialrules={\relentless{},\sturdy{}},
	weapons={\crossbow{}},
	armour={\ha{}},
	options={
		\mayswapcrossbowfor{
			\handgun{}=\free,
			\guildcraftedhandgun{} (\oneofakind{})=\permodel{}<2,
		},
		\shield{}=\permodel{}<1,
		\gw{}=\permodel{}<3,
	},
	unitequipment={
		\equipmentdef{\guildcraftedhandgun{}}{\guildcraftedhandgunrule{}}
	},
	commandgroup={champion=10, musician=10, banner=10, veteranstandardbearer=yessir},
}	











\specialunitstitle

\showunit{
	name={\deepwatch},
	QRSname={\deepwatchSINGULAR},
	cost={110},
	profile={ < 3 5 3 4 4 1 2 1 9},
	costpermodel=15,
	maxmodels=30,
	type=\infantry{},
	unitsize=10,
	basesize=20x20,
	commontype=\dwarvencommonrules{},
	commonspecialrules={\relentless{},\shieldwall{},\sturdy{}},
	specialrules={\bodyguard{},\shieldskill{}},
	armour={\platearmour{},\shield{}},
	unitrules={
		\unitrule{\shieldskill{}}{\shieldskillrule{}}
	},
	commandgroup={champion=10, musician=10, banner=10, bannerallowance=50},
}

\showunit{
	name={\kingsguard},
	cost={130},
	profile={ < 3 5 3 4 4 1 2 2 9},
	costpermodel=16,
	maxmodels=30,
	type=\infantry{},
	unitsize=10,
	basesize=20x20,
	commontype=\dwarvencommonrules{},
	commonspecialrules={\relentless{},\sturdy{}},
	specialrules={\bodyguard{\general{}, \king{}}},
	armour={\platearmour{}},
	weapons={\gw{}},
	options={
		\shield{}=\permodel{}<1,
	},
	commandgroup={champion=10, musician=10, banner=10, bannerallowance=50},
}

\showunit{
	name={\forgewardens},
	QRSname={\forgewarden},
	cost={140},
	profile={ < 3 5 3 4 4 1 2 1 9},
	costpermodel=14,
	maxmodels=20,
	type=\infantry{},
	unitsize=10,
	basesize=20x20,
	commontype=\dwarvencommonrules{},
	commonspecialrules={\relentless{},\sturdy{}},
	specialrules={\fireborn{},\wardsave{6}},
	armour={\platearmour{}},
	weapons={\forgegun{}},
	unitequipment={
		\equipmentdef{\forgegun{}}{\forgegunrule{}}
	},
	commandgroup={champion=10, championoption=\wyrmslayerrocket{}<20, musician=10, banner=10, bannerallowance=25},
}

\showunit{
	name={\seekers},
	QRSname={\seeker},
	cost={55},
	profile={ < 3 4 3 4 4 1 2 1 10},
	costpermodel=11,
	maxmodels=25,
	type=\infantry{},
	unitsize=5,
	basesize=20x20,
	commontype=\dwarvencommonrules{},
	commonspecialrules={\relentless{},\sturdy{}},
	commontypeB=\seekerscommonrules{},
	commonspecialrulesB={\unbreakable{},\weaponmaster{},\wardsave{6},\thebiggertheyare{},\comewithme{}},
	weapons={\gw{},\pw{}},
	options={
		\vanguard{}=\permodel{}<2,
		\seekersskirmishersoption{}=\permodel{}<2,
	},
	commandgroup={champion=10,musician=10},
}

\showunit{
	name={\steamcopters},
	QRSname={\steamcopter},
	profile={
		\copter{}< 1 - - - 5 3 - - -,
		\pilot{}< - 4 3 4 4 - 2 2 9,
		},
	type=\cavalry{},
	basesize=40x40,
	specialrules={\fly{8}},
	armour={\ha{},\mountsprotection{6}},
	additional={
		\begin{center}\copterchoice{}\end{center}
		\setlength{\columnseprule}{0.5pt}
		\renewcommand{\columnseprulecolor}{\color{black!30}}	
		\begin{multicols}{3}
		\raggedcolumns
		\begin{center}{\largerfontsize\antiquefont\attackcopter{} (\pts{75})}
		
		\vspace{5pt}\zerotoXchoice{2}\end{center}
		\textbf{1} \labels@model{}.\newline \mayaddonemodelfor{} \pts{\textbf{65}}.
		
		\vspace{0.15cm}
		\def\tempweapons{\forgerepeater{},\shrapnelgrenades{}}
		\weapons{\tempweapons}
		
		\def\tempoptions{%
			\copterskirmishoption{}=10,
		}
		\vspace*{0.1cm}\options{\tempoptions}

		\columnbreak
		\begin{center}{\largerfontsize\antiquefont\steambomber{} (\pts{105})}
		
		\vspace{5pt}\zerotoXchoice{2}\end{center}
		\labels@Singlemodel{}.
		
		\vspace{0.15cm}
		\def\tempweapons{\forgerepeater{},\shrapnelbombs{}}
		\weapons{\tempweapons}
		
		\vspace*{\fill}\columnbreak	
		\begin{center}{\largerfontsize\antiquefont\flamecopter{} (\pts{85})}
		
		\vspace{5pt}\zerotoXchoice{2}\end{center}
		\labels@Singlemodel{}.
		
		\vspace{0.15cm}
		\def\tempweapons{\flameculverin}
		\weapons{\tempweapons}	
		
		\vspace*{\fill}
		\end{multicols}
		\setlength{\columnseprule}{0pt}
		
		\def\tempunitequipment{%
			\equipmentdef{\shrapnelgrenades}{\shrapnelgrenadesrule}
			\equipmentdef{\shrapnelbombs}{\shrapnelbombsrule}
			\equipmentdef{\flameculverin}{\flameculverinrule}
		}
		\vspace*{0.1cm}\unitequipment{\tempunitequipment}
	},
}

\showunit{
	name={\miners},
	QRSname={\miner},
	cost={90},
	profile={ < 3 4 3 4 4 1 2 1 9},
	costpermodel=10,
	maxmodels=20,
	type=\infantry{},
	unitsize=10,
	basesize=20x20,
	commontype=\dwarvencommonrules{},
	commonspecialrules={\relentless{},\sturdy{}},
	specialrules={\ambush{}},
	armour={\ha{}},
	options={
		\musttakeoneormoreofthefollowing{
			\shield{}=\permodel{}<1,
			\throwingweapons{}=\permodel{}<2,
			\pistol{}=\permodel{}<2,	
			\gw{}=\permodel{}<2
		},
	},
	commandgroup={champion=10, musician=10, banner=10},
}

\showunit{
	name={\fieldartillery},
	profile={
		\machine{} < - - - - 7 3 - - -,
		\crew{} (3) < 3 4 3 3 4 - 2 1 9,
	},
	type=\warmachine{},
	unitsize=1,
	basesize=60,
	specialrules={\stubborn{}},
	armour={\ha{}},
	options={
		\flamingshot{}=5,
		\engineeringrune{}=15,
	},
	unitrules={
		\unitrule{\flamingshot}{\flamingshotrule}
		\unitrule{\engineeringrune}{\engineeringrunerule}
	},
	additional={
		\begin{center}\musttakeoneofthefollowingNOC{}\end{center}
		\setlength{\columnseprule}{0.5pt}
		\renewcommand{\columnseprulecolor}{\color{black!30}}	
		\begin{multicols}{3}
		\raggedcolumns		

		\begin{center}{\largerfontsize\antiquefont\catapult{} (\pts{90})}
		
		\vspace{5pt}\zerotoXchoice{2}\end{center}

		\noindent\catapultrule{}
		
		\def\tempunitrules{\unitrule{\runecrafted{} \textnormal{(\pts{50})}}{\catapultrunecraftedrule}}
		\vspace{0.2cm}\noindent\maytake{}\spacebeforecolon{}:
		\vspace*{0.1cm}\unitrules{\tempunitrules}

		\vspace*{\fill}\columnbreak
		\begin{center}{\largerfontsize\antiquefont\cannon{} (\pts{100})}
		
		\vspace{5pt}\zerotoXchoice{2}\end{center}
		
		\noindent\cannonrule{}
		
		\def\tempunitrules{\unitrule{\runecrafted{} \textnormal{(\pts{20})}}{\cannonrunecraftedrule}}
		\vspace{0.2cm}\noindent\maytake{}\spacebeforecolon{}:
		\vspace*{0.1cm}\unitrules{\tempunitrules}
		
		\columnbreak
		\begin{center}{\largerfontsize\antiquefont\organgun{} (\pts{125})}
		
		\vspace{5pt}\zerotoXchoice{2}\end{center}
		
		\noindent\organgunrule{}
		
		\def\tempunitrules{\unitrule{\runecrafted{} \textnormal{(\pts{45})}}{\organgunrunecraftedrule}}
		\vspace{0.2cm}\noindent\maytake{}\spacebeforecolon{}:
		\vspace*{0.1cm}\unitrules{\tempunitrules}

		\vspace*{\fill}
		\end{multicols}
		\setlength{\columnseprule}{0pt}
	},
}



	
	
	

\rareunitstitle

\showunit{
	name={\vengeanceseeker},
	cost={60},
	profile={< 3 5 3 4 4 3 2 \starsymbol{} 10},
	type=\infantry{},
	unitsize=1,
	basesize=20x20,
	commontype=\dwarvencommonrules{},
	commonspecialrules={\relentless{}},
	commontypeB=\seekerscommonrules{},
	commonspecialrulesB={\thebiggertheyare{},\comewithme{},\unbreakable{},\wardsave{6}},
	specialrules={\hardtarget{},\distracting{}},
	weapons={\whirlingchainsofdoom},
	unitequipment={
		\equipmentdef{\whirlingchainsofdoom{}}{\whirlingchainsofdoomrule{}}
	},
}
	
\showunit{
	name={\rangers},
	QRSname={\ranger},
	cost={65},
	profile={ < 3 4 4 3 4 1 2 1 9},
	costpermodel=10,
	maxmodels=20,
	type=\infantry{},
	unitsize=5,
	basesize=20x20,
	commontype=\dwarvencommonrules{},
	commonspecialrules={\relentless{},\sturdy{}},
	specialrules={\scouts{},\strider{\forest}},
	armour={\ha},
	options={
		\skirmishers{} \Xmodelsorless{10}=\permodel{}<2,
		\shield{}=\permodel{}<1,
		\throwingweapons{}=\permodel{}<1,
		\crossbow{}=\permodel{}<2,
		\combatweapononechoice{
			\pw{}=\permodel{}<1,
			\gw{}=\permodel{}<2,
		},
	},
	commandgroup={champion=10, musician=10, banner=10},
}				

\showunit{
	name={\holdguardians},
	QRSname={\holdguardian},
	cost={155},
	profile={< 5 4 3 6 5 3 2 2 10},
	costpermodel=60 ,
	maxmodels=8,
	type=\monstrousinfantry{},
	unitsize=3,
	basesize=40x40,
	specialrules={\magicalattacks{},\multiplewounds{1D3}{},\immunetopsychology{}},
	armour={\platearmour{},\innatedefence{6}},
	options={
		\flamingattacks{} \wordand{} \fireborn{}=\permodel{}<5,
	},
	commandgroup={champion=10, musician=10, banner=10},
}	

\showunit{
	name={\grudgebuster},
	cost={165},
	profile={
		\chariot{}< 1 - - 5 5 5 2 - -,
		\grudgebustercrew{} (2)< - 4 3 4 4 - 2 2 9,
	},
	type=\chariot{},
	unitsize=1,
	basesize=50x100,
	specialrules={\grindingattacks{3D3} \only{\chariot},\impacthits{+1},\fly{8}},
	weapons={\swiwelgun{} \only{\chariot}},
	armour={\innatedefence{3}},
	unitequipment={
		\equipmentdef{\swiwelgun{}}{\swiwelgunrule{}}
	},
}		

\showunit{
	name={\grudgemachine},
	profile={
		\machine{}< - - - - 7 3 - - -,
		\crew{} (3)< 3 4 3 3 4 - 2 1 9,
	},
	type=\warmachine{},
	unitsize=SPECIAL-{\textbf{(\oneofakind{})}\refsymbol{} \labels@Singlemodel},
	basesize=60,
	specialrules={\stubborn},
	armour={\ha},
	options={
		\engineeringrune{}=15,
	},
	unitrules={
		\unitrule{\engineeringrune}{\engineeringrunerule}
	},
	additional={
		\begin{center}\musttakeoneofthefollowingNOC{}\end{center}
		\setlength{\columnseprule}{0.5pt}
		\renewcommand{\columnseprulecolor}{\color{black!30}}	
		\begin{multicols}{2}
		\raggedcolumns		

		\begin{center}\largerfontsize\antiquefont\flamecannon{} (\pts{105})\end{center}

		\noindent\flamecannonrule{}
		
		\def\tempunitrules{%
			\unitrule{\runecrafted{} \textnormal{(\pts{5})}}{\flamecannonrunecraftedrule}
			\unitrule{\runestruck{} \textnormal{(\pts{20})}}{\runestruckrule}
		}
		\vspace{0.2cm}\noindent\maytake{}\spacebeforecolon{}:
		\vspace*{0.1cm}\unitrules{\tempunitrules}

		\vspace*{\fill}\columnbreak
		\begin{center}\largerfontsize\antiquefont\dwarfballista{} (\pts{55})\end{center}
		
		\noindent\dwarfballistarule{}
		
		\def\tempunitrules{%
			\unitrule{\flamingshot{} \textnormal{(\pts{5})}}{\flamingshotrule}
			\unitrule{\runecrafted{} \textnormal{(\pts{10})}}{\dwarfballistarunecraftedrule}
		}
		\vspace{0.2cm}\noindent\maytake{}\spacebeforecolon{}:
		\vspace*{0.1cm}\unitrules{\tempunitrules}
		\vspace*{\fill}
		\end{multicols}
		\setlength{\columnseprule}{0pt}
		
		\vspace*{0.15cm}\noindent\refsymbol{} \grudgemachinenote{}
	},		
}	

\showunit{
	name={\anvilofpower},
	cost=150,
	profile={
		\anvil{}< - - - - 7 3 - - -,
		\anvilcrew{} (3)< 3 5 3 4 4 - 2 1 9,
	},
	type=\warmachine{},
	unitsize=SPECIAL-{\textbf{(\oneofakind)} \labels@Singlemodel},
	basesize=60,
	specialrules={\channel{},\runicanvil{},\unbreakable{},\magicresistance{2},\wardsave{5}},
	armour={\platearmour},
	additional={%
		\def\tempunitrules{\unitrule{\runicanvil}{\runicanvilrule}}
		\unitrules{\tempunitrules}
	
		\anvilrunes{}
	},			
}	






%%% Quick Reference Sheet - AB_qrs.tex is automatic and shouldn't be edited %%%

\quickrefsheettitle

% Script to automatically draw the Quick Ref Sheet

\renewcommand{\arraystretch}{1.2}

\providebool{QRSbool}

\providebool{whiterow}

\newcommand{\QRSrowcolor}{\ifbool{whiterow}{\global\boolfalse{whiterow}}{\rowcolor{black!10}\global\booltrue{whiterow}}}

\newcommand{\QRSstarttab}[1]{%
	\noindent%
	\setlength{\tabcolsep}{2pt}%
	\begin{tabular}{@{}cp{3.2cm}M{\profilecellsize}@{}M{\profilecellsize}@{}M{\profilecellsize}@{}M{\profilecellsize}@{}M{\profilecellsize}@{}M{\profilecellsize}@{}M{\profilecellsize}@{}M{\profilecellsize}@{}M{\profilecellsize}}%

	& \antiquefont\large{\textbf{#1}} & \textbf{\labels@M} & \textbf{\labels@WS} & \textbf{\labels@BS} & \textbf{\labels@S} & \textbf{\labels@T} & \textbf{\labels@W} & \textbf{\labels@I} & \textbf{\labels@A} & \textbf{\labels@Ld}%
}%

\newcommand{\QRSclosetab}{\end{tabular}\bigskip}%

\newcommand{\QRSprintline}[4]{%
	\tabularnewline%
	\ifnumequal{\rowmulti}{1}{\QRSrowcolor}{}%
	\DTLifeq*{\rowcategory}{\labels@lords}{\antiquefont\bfseries \labels@lordsInitial}{}%
	\DTLifeq*{\rowcategory}{\labels@heroes}{\antiquefont\bfseries \labels@heroesInitial}{}%
	\DTLifeq*{\rowcategory}{\labels@coreunits}{\antiquefont\bfseries \labels@coreunitsInitial}{}%
	\DTLifeq*{\rowcategory}{\labels@specialunits}{\antiquefont\bfseries \labels@specialunitsInitial}{}%
	\DTLifeq*{\rowcategory}{\labels@rareunits}{\antiquefont\bfseries \labels@rareunitsInitial}{}%
	\DTLifeq*{\rowcategory}{\labels@mounts}{\antiquefont\bfseries \labels@mountsInitial}{}%
	&%
	\ifnumequal{\rowmulti}{1}{%no Multiprofile
		\rowname%
		\expandafter\parselist\expandafter{\rowprofile}{\locallists@profileslist}%
		\forlistloop{\QRSmonoprofile}{\locallists@profileslist}%
	}{% Multiprofile
		\rowname &&&&&&&&&%
		\expandafter\parselist\expandafter{\rowprofile}{\locallists@profileslist}%
		\forlistloop{\QRSmultiprofile}{\locallists@profileslist}%
	}%
}

\newcommand{\QRSmultiprofile}[1]{%
	\tabularnewline%
	\QRSrowcolor{}&%
	\splitatinf{#1}\local@unitname\local@unitprofile%
	- \local@unitname \expandafter\caraclist\expandafter{\local@unitprofile}%
}%

\newcommand{\QRSmonoprofile}[1]{%
	\splitatinf{#1}\local@unitname\local@unitprofile%
	\expandafter\caraclist\expandafter{\local@unitprofile}%
}%

\newcommand{\QRSprinttab}[1]{%
	\global\booltrue{whiterow}%
	\DTLforeach*[#1]%
	{profiles}{\rowname=name, \rowtrooptype=trooptype, \rowcategory=category, \rowprofile=profile, \rowmulti=multipleprofile}{%
      		\QRSprintline{\rowname}{\rowcategory}{\rowprofile}{\rowmulti}%
	}%
}%

\providebool{QRSisempty}
\global\boolfalse{QRSisempty}%

\newcommand{\QRScheckifempty}[1]{%
	\global\booltrue{QRSisempty}%
	\DTLforeach*[#1]%
	{profiles}{\rowname=name, \rowtrooptype=trooptype, \rowcategory=category, \rowprofile=profile, \rowmulti=multipleprofile}{%
		\global\boolfalse{QRSisempty}\dtlbreak%
	}%
}%

\newcommand{\QRSifnotempty}[1]{%
	\ifbool{QRSisempty}{}{#1}%
}%

\begin{center}
{\antiquefont\bfseries \labels@lordsInitial}\spacebeforecolon{}: \labels@lords{} - %
{\antiquefont\bfseries \labels@heroesInitial}\spacebeforecolon{}: \labels@heroes{} - %
{\antiquefont\bfseries \labels@coreunitsInitial}\spacebeforecolon{}: \labels@coreunits{} - %
{\antiquefont\bfseries \labels@specialunitsInitial}\spacebeforecolon{}: \labels@specialunits{} - %
{\antiquefont\bfseries \labels@rareunitsInitial}\spacebeforecolon{}: \labels@rareunits{} - %
{\antiquefont\bfseries \labels@mountsInitial}\spacebeforecolon{}: \labels@mounts{}%
\end{center}

\begin{multicols}{2}

\QRScheckifempty{%
	\DTLiseq{\rowcategory}{\labels@lords}\or\DTLiseq{\rowcategory}{\labels@heroes}%
}%
\QRSifnotempty{%
	\QRSstarttab{\characters}%
	\QRSprinttab{%
		\DTLiseq{\rowcategory}{\labels@lords}\or\DTLiseq{\rowcategory}{\labels@heroes}%
	}%
	\QRSclosetab{}%
}%

\QRScheckifempty{%
	\DTLiseq{\rowtrooptype}{\infantry}\and\not\DTLiseq{\rowcategory}{\labels@heroes}\and\not\DTLiseq{\rowcategory}{\labels@lords}%
}%
\QRSifnotempty{%
	\QRSstarttab{\infantry}%
	\QRSprinttab{%
		\DTLiseq{\rowtrooptype}{\infantry}\and\not\DTLiseq{\rowcategory}{\labels@heroes}\and\not\DTLiseq{\rowcategory}{\labels@lords}%
	}% 
	\QRSclosetab{}%
}% 

\QRScheckifempty{%
	\DTLiseq{\rowtrooptype}{\monstrousinfantry}\and\not\DTLiseq{\rowcategory}{\labels@heroes}\and\not\DTLiseq{\rowcategory}{\labels@lords}%
}%
\QRSifnotempty{%
	\QRSstarttab{\monstrousinfantry}%
	\QRSprinttab{%
		\DTLiseq{\rowtrooptype}{\monstrousinfantry}\and\not\DTLiseq{\rowcategory}{\labels@heroes}\and\not\DTLiseq{\rowcategory}{\labels@lords}%
	}% 
	\QRSclosetab{}%
}% 

\QRScheckifempty{%
	\DTLiseq{\rowtrooptype}{\warbeast}\and\not\DTLiseq{\rowcategory}{\labels@heroes}\and\not\DTLiseq{\rowcategory}{\labels@lords}%
}%
\QRSifnotempty{%
	\QRSstarttab{\warbeasts}%
	\QRSprinttab{%
		\DTLiseq{\rowtrooptype}{\warbeast}\and\not\DTLiseq{\rowcategory}{\labels@heroes}\and\not\DTLiseq{\rowcategory}{\labels@lords}%
	}% 
	\QRSclosetab{}%
}% 

\QRScheckifempty{%
	\DTLiseq{\rowtrooptype}{\monstrousbeast}\and\not\DTLiseq{\rowcategory}{\labels@heroes}\and\not\DTLiseq{\rowcategory}{\labels@lords}%
}%
\QRSifnotempty{%
	\QRSstarttab{\monstrousbeasts}%
	\QRSprinttab{%
		\DTLiseq{\rowtrooptype}{\monstrousbeast}\and\not\DTLiseq{\rowcategory}{\labels@heroes}\and\not\DTLiseq{\rowcategory}{\labels@lords}%
	}% 
	\QRSclosetab{}%
}% 

\QRScheckifempty{%
	\DTLiseq{\rowtrooptype}{\cavalry}\and\not\DTLiseq{\rowcategory}{\labels@heroes}\and\not\DTLiseq{\rowcategory}{\labels@lords}%
}%
\QRSifnotempty{%
	\QRSstarttab{\cavalry}%
	\QRSprinttab{%
		\DTLiseq{\rowtrooptype}{\cavalry}\and\not\DTLiseq{\rowcategory}{\labels@heroes}\and\not\DTLiseq{\rowcategory}{\labels@lords}%
	}%
	\QRSclosetab{}%
}% 

\QRScheckifempty{%
	\DTLiseq{\rowtrooptype}{\monstrouscavalry}\and\not\DTLiseq{\rowcategory}{\labels@heroes}\and\not\DTLiseq{\rowcategory}{\labels@lords}%
}%
\QRSifnotempty{%
	\QRSstarttab{\monstrouscavalry}%
	\QRSprinttab{%
		\DTLiseq{\rowtrooptype}{\monstrouscavalry}\and\not\DTLiseq{\rowcategory}{\labels@heroes}\and\not\DTLiseq{\rowcategory}{\labels@lords}%
	}%
	\QRSclosetab{}%
}% 

\QRScheckifempty{%
	\DTLiseq{\rowtrooptype}{\chariot}\and\not\DTLiseq{\rowcategory}{\labels@heroes}\and\not\DTLiseq{\rowcategory}{\labels@lords}%
}%
\QRSifnotempty{%
	\QRSstarttab{\chariots}%
	\QRSprinttab{%
		\DTLiseq{\rowtrooptype}{\chariot}\and\not\DTLiseq{\rowcategory}{\labels@heroes}\and\not\DTLiseq{\rowcategory}{\labels@lords}%
	}%
	\QRSclosetab{}%
}% 

\QRScheckifempty{%
	\DTLiseq{\rowtrooptype}{\monster}\and\not\DTLiseq{\rowcategory}{\labels@heroes}\and\not\DTLiseq{\rowcategory}{\labels@lords}%
}%
\QRSifnotempty{%
	\QRSstarttab{\monsters}%
	\QRSprinttab{%
		\DTLiseq{\rowtrooptype}{\monster}\and\not\DTLiseq{\rowcategory}{\labels@heroes}\and\not\DTLiseq{\rowcategory}{\labels@lords}%
	}%
	\QRSclosetab{}%
}% 

\QRScheckifempty{%
	\DTLiseq{\rowtrooptype}{\riddenmonster}\and\not\DTLiseq{\rowcategory}{\labels@heroes}\and\not\DTLiseq{\rowcategory}{\labels@lords}%
}%
\QRSifnotempty{%
	\QRSstarttab{\riddenmonsters}%
	\QRSprinttab{%
		\DTLiseq{\rowtrooptype}{\riddenmonster}\and\not\DTLiseq{\rowcategory}{\labels@heroes}\and\not\DTLiseq{\rowcategory}{\labels@lords}%
	}%
	\QRSclosetab{}%
}% 

\QRScheckifempty{%
	\DTLiseq{\rowtrooptype}{\swarm}\and\not\DTLiseq{\rowcategory}{\labels@heroes}\and\not\DTLiseq{\rowcategory}{\labels@lords}%
}%
\QRSifnotempty{%
	\QRSstarttab{\swarms}%
	\QRSprinttab{%
		\DTLiseq{\rowtrooptype}{\swarm}\and\not\DTLiseq{\rowcategory}{\labels@heroes}\and\not\DTLiseq{\rowcategory}{\labels@lords}%
	}%
	\QRSclosetab{}%
}% 

\end{multicols}
\bigskip
\begin{center}
\noindent\begin{center}\Largefontsize{\textbf{Armes de Tir Naines}}\end{center}
\medskip

\rowcolors{1}{white}{black!10}
\noindent\begin{tabular}{lcccccc}
\textbf{Nom} & \textbf{Artillerie} & \textbf{Portée} & \textbf{\labels@S{}} & \textbf{\multipleshots{}} & \textbf{\multiplewounds{}} & \textbf{\armourpiercing{}} \tabularnewline
\forgegun{} (\forgewardens{}) & - & \distance{18} & 5 & - & - & - \tabularnewline
\forgerepeater{} & - & \distance{18} & 5 & 4 & - & - \tabularnewline
\wyrmslayerrocket{} & - & \distance{24} & 6 & - & 1D3 & - \tabularnewline
\swiwelgun{} (\grudgebuster{}) & - & \distance{24} & 4 & 4 & - & 1 \tabularnewline
\flameculverin{} (\copter{}) & \flamethrower{} & \distance{3} & 3 & - & - & -\tabularnewline
\fieldartillery{} - \cannon{} (1) & \cannon{} (\distance{1D6}) & \distance{60} & 10 & - & \ordnance{} & 2 \tabularnewline
\fieldartillery{} - \cannon{} (2) & \volleygun{} & \distance{12} & 4 & 2D6 & - & 3 \tabularnewline
\fieldartillery{} - \organgun{} & \volleygun{} & \distance{30} & 5 & 2D6*2 & - & 1 \tabularnewline
\fieldartillery{} - \catapult{} & \catapult{} (\distance{3}) & \distance{60} & 3 [9] & - & [\ordnance{}] & - \tabularnewline
\grudgemachine{} - \dwarfballista{} & \boltthrower{} & \distance{48} & 6 & - & 1D3 & 6 \tabularnewline
\grudgemachine{} - \flamecannon{} & \flamethrower{} & \distance{12} & 5 & - & 1D3 & - \tabularnewline
\end{tabular}

\end{center}

\restoregeometry

\end{document}